%% ============================================================================
%%  CAPITOLO 7 — Funzioni di più variabili
%%  Fonte principale per argomenti e dimostrazioni: Marcellini–Sbordone, cap. 7
%%  (§54 funzioni di due variabili, dominio, rappresentazione cartesiana;
%%  §55 limiti e continuità; §56 derivate parziali, gradiente; §57 derivate
%%  successive, teorema di Schwarz; §58 massimi e minimi relativi; §59 funzioni
%%  di tre o più variabili; appendice §60 differenziabilità).
%% ============================================================================
\aperturacapitolo{Funzioni di più variabili}{cap:piu-variabili}

\epigraph{On ne trouvera point de Figures dans cet Ouvrage. Les méthodes que j'y
expose ne demandent ni constructions, ni raisonnemens géométriques ou méchaniques, mais
seulement des opérations algébriques.\\[2pt]
\upshape(Non si troveranno figure in quest'opera. I metodi che vi espongo non
richiedono né costruzioni, né ragionamenti geometrici o meccanici, ma soltanto
operazioni algebriche.)}{Joseph-Louis Lagrange, \emph{Méchanique analitique} (1788),
avvertenza}

\lettrine[lines=3, findent=2pt, nindent=0pt]{Q}{uasi tutte} le grandezze che un
ingegnere incontra dipendono da più di una variabile. La portanza di un'ala dipende
dalla densità dell'aria, dalla velocità e dall'assetto; la pressione atmosferica dipende
dalla posizione e dalla quota; la sollecitazione in un longherone dipende dal punto
della sezione e dalla posizione lungo l'apertura. In questo capitolo diamo alcuni cenni
dell'analisi delle funzioni di più variabili reali, un argomento che in genere viene
ripreso e approfondito in un corso di Analisi Matematica di secondo anno, ma che è
troppo importante per le applicazioni per rimandarlo del tutto.

Vedremo che molte idee del calcolo per le funzioni di una variabile si estendono in modo
naturale: il limite e la continuità si definiscono sostituendo la distanza nel piano al
valore assoluto; la derivata diventa la \emph{derivata parziale}, che si calcola
derivando rispetto a una variabile e tenendo ferme le altre; la condizione di Fermat
$f'(x_{0})=0$ diventa l'annullarsi del \emph{gradiente}. Vedremo anche che non tutto si
estende senza sorprese: una funzione può avere entrambe le derivate parziali in un punto
senza essere continua in quel punto, e le derivate seconde miste possono dipendere
dall'ordine di derivazione. Queste sorprese ci condurranno, in appendice, alla nozione
di \emph{differenziabilità}, la vera generalizzazione della derivabilità. Malgrado
l'epigrafe di Lagrange, il capitolo è ricchissimo di figure: per le funzioni di due
variabili il grafico è una superficie, e vederla aiuta moltissimo. Seguiremo il testo di
riferimento \cite{marcellini-sbordone}.

\medskip
\noindent\textbf{In questo capitolo.}
\begin{itemize}
  \item Funzioni di due variabili, grafici, curve di livello, domini
  (\S\ref{sec:due-variabili}).
  \item Limiti e continuità nel piano (\S\ref{sec:limiti-piano}).
  \item Derivate parziali, gradiente, derivate direzionali (\S\ref{sec:derivate-parziali}).
  \item Derivate successive e teorema di Schwarz; l'equazione di Laplace
  (\S\ref{sec:schwarz}).
  \item Approfondimento: un primo sguardo alle equazioni alle derivate parziali
  (\S\ref{subsec:edp-primo-sguardo}).
  \item Massimi e minimi relativi; il metodo dei minimi quadrati (\S\ref{sec:estremi-due}).
  \item Funzioni di tre o più variabili (\S\ref{sec:n-variabili}).
  \item Appendice: differenziabilità, piano tangente, dimostrazione del criterio
  dell'hessiano (\S\ref{sec:differenziabilita}).
  \item Approfondimento: funzioni omogenee e teorema di Eulero
  (\S\ref{subsec:funzioni-omogenee}).
  \item Esercizi di riepilogo (\S\ref{sec:esercizi-cap7}).
  \item Punti fondamentali del capitolo (\S\ref{sec:sintesi-cap07}).
\end{itemize}

%% ----------------------------------------------------------------------------
%%  §7.1  Funzioni di due variabili: dominio, grafico, curve di livello
%% ----------------------------------------------------------------------------
\section{Funzioni di due variabili: dominio, grafico, curve di livello}
\label{sec:due-variabili}

Indichiamo con $\R^{2}$ l'insieme delle coppie ordinate di numeri reali:
\begin{equation}\label{eq:R2}
  \R^{2}=\set{(x,y) : x\in\R,\ y\in\R}.
\end{equation}
Fissato nel piano un riferimento cartesiano ortogonale, ogni coppia $(x,y)$ individua un
punto $P$ del piano, e viceversa. La \emph{distanza} tra due punti $P_{1}=(x_{1},y_{1})$ e
$P_{2}=(x_{2},y_{2})$ è, per il teorema di Pitagora,
\begin{equation}\label{eq:distanza}
  d(P_{1},P_{2})=\sqrt{(x_{1}-x_{2})^{2}+(y_{1}-y_{2})^{2}} .
\end{equation}
Ci servirà spesso il confronto tra la distanza e le differenze delle coordinate: posto
$a=x_{1}-x_{2}$ e $b=y_{1}-y_{2}$, poiché $a^{2}\le a^{2}+b^{2}\le(\abs a+\abs b)^{2}$,
risulta
\begin{equation}\label{eq:distanza-componenti}
\begin{gathered}
  \abs{x_{1}-x_{2}}\le d(P_{1},P_{2}), \qquad \abs{y_{1}-y_{2}}\le d(P_{1},P_{2}),\\
  d(P_{1},P_{2})\le\abs{x_{1}-x_{2}}+\abs{y_{1}-y_{2}} .
\end{gathered}
\end{equation}
In parole: due punti sono vicini se e soltanto se sono vicine sia le ascisse sia le
ordinate.

\begin{definizione}[Funzione di due variabili]\label{def:funzione-due-variabili}
\index{funzione!di due variabili}\index{dominio}
Sia $D$ un sottoinsieme di $\R^{2}$. Una legge $f$ che a ogni elemento di $D$ fa
corrispondere uno e un solo numero reale si dice una \emph{funzione di due variabili}; si
scrive $f:D\to\R$, oppure, per mettere in evidenza che $D\subseteq\R^{2}$,
\begin{equation}\label{eq:z-f-xy}
  z=f(x,y), \qquad (x,y)\in D .
\end{equation}
L'insieme $D$ si dice il \emph{dominio} (o \emph{insieme di definizione}) di $f$.
L'insieme dei punti $(x,y,z)$ dello spazio con $(x,y)\in D$ e $z=f(x,y)$ si chiama il
\emph{grafico} di $f$: è, in generale, una superficie.
\end{definizione}

Per rappresentare graficamente una funzione di due variabili si usa un riferimento
cartesiano ortogonale di assi $x$, $y$, $z$: a ogni punto $(x,y)$ del dominio, nel piano
di base, si associa il punto $(x,y,z)$ a quota $z=f(x,y)$. Un ingegnere aeronautico
incontra funzioni di due variabili a ogni passo. La portanza di un velivolo al livello
del mare, per una data superficie alare $S$, è
\begin{equation}\label{eq:portanza-due-variabili}
  L(V,C_{L})=\frac12\rho S V^{2}C_{L},
\end{equation}
una funzione della velocità $V$ e del coefficiente di portanza $C_{L}$. Prima ancora dei
grafici, il modo tradizionale di rappresentare una funzione di due variabili è la
\emph{tabella a doppia entrata}, come quelle dell'atmosfera standard o dei prontuari.

\begin{python}[Una tabella a doppia entrata: la portanza $L(V,C_{L})$]
rho, S, W = 1.225, 16.0, 10000.0         # kg/m^3, m^2, N
L = lambda V, CL: 0.5 * rho * S * V**2 * CL

CLs = (0.2, 0.4, 0.6, 0.8, 1.0, 1.2)
print(" V (m/s) |" + "".join(f"   CL={c:.1f}" for c in CLs))
for V in (30, 40, 50, 60, 70):
    riga = "".join(f"{L(V, c):9.0f}" for c in CLs)
    print(f"  {V:5d}  |{riga}")
# le coppie (V, C_L) per cui L = W: C_L = 2W / (rho S V^2)
for V in (30, 40, 50, 60, 70):
    print(f"V = {V} m/s:  L = W per C_L = "
          f"{2 * W / (rho * S * V**2):.3f}")
\end{python}
\begin{risultato}
 V (m/s) |   CL=0.2   CL=0.4   CL=0.6   CL=0.8   CL=1.0   CL=1.2
     30  |     1764     3528     5292     7056     8820    10584
     40  |     3136     6272     9408    12544    15680    18816
     50  |     4900     9800    14700    19600    24500    29400
     60  |     7056    14112    21168    28224    35280    42336
     70  |     9604    19208    28812    38416    48020    57624
V = 30 m/s:  L = W per C_L = 1.134
V = 40 m/s:  L = W per C_L = 0.638
V = 50 m/s:  L = W per C_L = 0.408
V = 60 m/s:  L = W per C_L = 0.283
V = 70 m/s:  L = W per C_L = 0.208
\end{risultato}

Le ultime righe rispondono alla domanda che interessa al pilota: per sostenere il peso
$W=10$~kN, quale coefficiente di portanza serve a ciascuna velocità? Le coppie $(V,C_{L})$
per cui $L(V,C_{L})=W$ formano una curva del piano $(V,C_{L})$, che è un esempio della
nozione seguente.

\begin{definizione}[Curve di livello]\label{def:curve-livello}\index{curva di livello}
Sia $f:D\to\R$ e sia $c\in\R$. L'insieme
\begin{equation}\label{eq:curva-livello}
  \set{(x,y)\in D : f(x,y)=c}
\end{equation}
si chiama \emph{insieme di livello} (o, quando è una curva, \emph{curva di livello}) di
$f$ relativo al valore $c$.
\end{definizione}

La curva di livello $c$ è la proiezione sul piano di base della sezione del grafico con
il piano orizzontale $z=c$. Le carte topografiche rappresentano il terreno con le sue
curve di livello (le isoipse), le carte del tempo rappresentano la pressione con le
isobare: in entrambi i casi una funzione di due variabili è descritta, su un foglio
piano, dalla famiglia delle sue curve di livello (\cref{fig:livello}). Dove le curve di
livello sono fitte, la funzione varia rapidamente; dove sono rade, varia lentamente.

\begin{figure}[htb]
\centering
\includegraphics[width=\textwidth]{livello}
\caption{Due funzioni di due variabili descritte con le loro curve di livello. A
sinistra, le isobare (in hPa) di un campo di pressione con un'alta $A$ e una bassa $B$; le
frecce sono i vettori $-\nabla p$, di cui parleremo nel \cref{sec:derivate-parziali}. A
destra, le curve di livello della portanza $L(V,C_{L})$ di un velivolo con $S=16$~m$^{2}$
al livello del mare; la curva in rame è quella del volo livellato, $L=W$
(\texttt{python/fig\_livello.py}).}
\label{fig:livello}
\end{figure}

\subsection{Esempi di grafici}

Il comportamento di una funzione di due variabili si comprende spesso \emph{fissando una
delle due variabili} e studiando la funzione di una variabile che si ottiene.
Consideriamo
\begin{equation}\label{eq:sella}
  f(x,y)=x^{2}-y^{2},
\end{equation}
definita per ogni $(x,y)\in\R^{2}$. Per $y$ fissato si ottengono le parabole
\emph{convesse} $z=x^{2}-\text{costante}$, per $x$ fissato le parabole \emph{concave}
$z=\text{costante}-y^{2}$. Il grafico, rappresentato nella \cref{fig:superfici} sul
quadrato $[-1,1]\times[-1,1]$, ha la forma di una sella e si chiama \emph{paraboloide
iperbolico}: le sue curve di livello $x^{2}-y^{2}=c$ sono iperboli. Scambiando i ruoli
delle variabili si ottiene il grafico di $f(x,y)=y^{2}-x^{2}$, ruotato di un angolo retto.

La funzione
\begin{equation}\label{eq:rigata}
  f(x,y)=y\,(x^{2}+x)
\end{equation}
ha, per $y>0$ fissato, l'andamento di una parabola convessa del tipo $z=x^{2}+x$, mentre
per $y<0$ fissato l'andamento è quello di una parabola concava. Per $x$ fissato, invece,
la funzione è lineare: si tratta della retta $z=\text{costante}\cdot y$. Il grafico è
quindi unione di una famiglia di rette, e per questo si dice una \emph{superficie
rigata}.

\begin{figure}[htb]
\centering
\includegraphics[width=0.95\textwidth]{superfici}
\caption{I grafici di quattro funzioni di due variabili: il paraboloide iperbolico, la
superficie rigata $y(x^{2}+x)$, la funzione $\cos(x^{2}+y^{2})$, invariante per
rotazioni, e la funzione $\sen(xy)$, costante sulle iperboli $xy=t$
(\texttt{python/fig\_superfici.py}).}
\label{fig:superfici}
\end{figure}

La funzione
\begin{equation}\label{eq:cos-radiale}
  f(x,y)=\cos(x^{2}+y^{2})
\end{equation}
è definita per ogni $(x,y)\in\R^{2}$. Per comprenderne il grafico conviene pensarla come
funzione composta: posto $t=\sqrt{x^{2}+y^{2}}\ge0$, distanza del punto $(x,y)$
dall'origine, è $f(x,y)=\cos t^{2}$. La funzione è costante su ogni circonferenza
$x^{2}+y^{2}=t^{2}$ di centro l'origine, e il suo valore è $\cos t^{2}$. Il grafico si
ottiene quindi facendo ruotare intorno all'asse $z$ il profilo $z=\cos t^{2}$, $t\ge0$:
si dice che $f$ è \emph{invariante per rotazioni}. Più in generale, ogni funzione della
forma $f(x,y)=g\bigl(\sqrt{x^{2}+y^{2}}\bigr)$ è invariante per rotazioni, e il suo
grafico è una superficie di rotazione. La funzione
\begin{equation}\label{eq:sen-xy}
  f(x,y)=\sen(xy)
\end{equation}
è invece costante sulle iperboli $xy=t$, con $t$ fissato in $\R$.

Sono invarianti per rotazioni anche le funzioni
\begin{equation}\label{eq:radiali}
\begin{gathered}
  \log(x^{2}+y^{2}), \qquad f_{0}(x,y)=1-\sqrt{x^{2}+y^{2}},\\
  f_{1}(x,y)=1-\sqrt{x^{2}+y^{2}-1}, \qquad f_{2}(x,y)=1-\sqrt{x^{2}+y^{2}-2},
\end{gathered}
\end{equation}
i cui grafici sono nella \cref{fig:radiali}. La prima non è definita nell'origine, dove
$x^{2}+y^{2}=0$; il suo profilo è $z=\log t^{2}=2\log t$, $t>0$. La funzione $f_{0}$,
definita su tutto $\R^{2}$, ha per $y=0$ il profilo $z=1-\sqrt{x^{2}}=1-\abs x$: il grafico è
un cono, con il vertice appuntito nel punto $(0,0,1)$, dovuto al valore assoluto. La funzione
$f_{1}$ è definita se $x^{2}+y^{2}\ge1$, cioè all'esterno del cerchio di centro
l'origine e raggio $1$ (circonferenza compresa); analogamente $f_{2}$ è definita
all'esterno del cerchio di raggio $\sqrt2$.

\begin{figure}[htb]
\centering
\includegraphics[width=0.95\textwidth]{radiali}
\caption{Quattro funzioni invarianti per rotazioni: i grafici sono superfici di
rotazione intorno all'asse $z$. Le ultime due sono definite soltanto all'esterno di un
cerchio, sulla cui circonferenza il grafico ha tangente verticale
(\texttt{python/fig\_radiali.py}).}
\label{fig:radiali}
\end{figure}

\subsection{La determinazione del dominio}

Quando una funzione è assegnata mediante un'espressione analitica, senza specificare il
dominio, si sottintende che il dominio sia l'insieme di tutti i punti in cui
l'espressione ha significato. Determinarlo richiede di risolvere disequazioni in due
variabili, la cui soluzione è una regione del piano. Ad esempio, le funzioni
\begin{equation}\label{eq:domini-esempi}
\begin{gathered}
  z=1-\sqrt{(1-x^{2})(1-y^{2})},\\
  z=\sqrt{(x^{2}+y^{2}-x)(2x-x^{2}-y^{2})}
\end{gathered}
\end{equation}
non sono definite su tutto $\R^{2}$. La prima è definita quando
$(1-x^{2})(1-y^{2})\ge0$, cioè quando i fattori $1-x^{2}$ e $1-y^{2}$ hanno lo stesso
segno (o uno dei due è nullo): o entrambi $\ge0$, cioè $\abs x\le1$ e $\abs y\le1$ (il
quadrato centrale), oppure entrambi $\le0$, cioè $\abs x\ge1$ e $\abs y\ge1$ (le quattro
regioni d'angolo); si veda la \cref{fig:domini}, a sinistra.

La seconda è definita quando l'argomento della radice è $\ge0$. Completando i quadrati,
\begin{equation}\label{eq:due-cerchi}
\begin{gathered}
  x^{2}+y^{2}-x=\Bigl(x-\frac12\Bigr)^{2}+y^{2}-\frac14,\\
  2x-x^{2}-y^{2}=1-\bigl[(x-1)^{2}+y^{2}\bigr];
\end{gathered}
\end{equation}
il primo fattore è $\ge0$ all'esterno della circonferenza $C_{1}$ di centro
$(\frac12,0)$ e raggio $\frac12$, il secondo è $\ge0$ all'interno della circonferenza
$C_{2}$ di centro $(1,0)$ e raggio $1$ (circonferenze comprese). Il prodotto è $\ge0$
quando i fattori hanno lo stesso segno. Il caso «entrambi negativi» è impossibile: un
punto interno a $C_{1}$ è anche interno a $C_{2}$, perché $C_{1}$ è contenuta nel cerchio
delimitato da $C_{2}$ (le due circonferenze sono tangenti internamente nell'origine).
Resta il caso «entrambi $\ge0$»: il dominio è la lunula compresa tra le due
circonferenze (\cref{fig:domini}, a destra).

\begin{figure}[htb]
\centering
\begin{tikzpicture}[scale=0.85]
  \begin{scope}
    \clip (-2.3,-2.3) rectangle (2.3,2.3);
    \fill[pattern=north east lines, pattern color=blu] (-1,-1) rectangle (1,1);
    \fill[pattern=north east lines, pattern color=blu] (1,1) rectangle (2.3,2.3);
    \fill[pattern=north east lines, pattern color=blu] (-1,1) rectangle (-2.3,2.3);
    \fill[pattern=north east lines, pattern color=blu] (1,-1) rectangle (2.3,-2.3);
    \fill[pattern=north east lines, pattern color=blu] (-1,-1) rectangle (-2.3,-2.3);
    \draw[blu, thick] (-1,-2.3) -- (-1,2.3) (1,-2.3) -- (1,2.3)
      (-2.3,-1) -- (2.3,-1) (-2.3,1) -- (2.3,1);
  \end{scope}
  \draw[-Latex, grigio] (-2.5,0) -- (2.6,0) node[right, black] {$x$};
  \draw[-Latex, grigio] (0,-2.5) -- (0,2.6) node[above, black] {$y$};
  \node[font=\scriptsize, below left] at (-1,0) {$-1$};
  \node[font=\scriptsize, below right] at (1,0) {$1$};
  \node[font=\scriptsize, above left] at (0,1) {$1$};
  \node[font=\scriptsize, below left] at (0,-1) {$-1$};
\end{tikzpicture}\hspace{1cm}
\begin{tikzpicture}[scale=1.2]
  \fill[pattern=north east lines, pattern color=blu] (1,0) circle (1);
  \fill[white] (0.5,0) circle (0.5);
  \draw[blu, thick] (1,0) circle (1);
  \draw[blu, thick] (0.5,0) circle (0.5);
  \draw[-Latex, grigio] (-0.4,0) -- (2.4,0) node[right, black] {$x$};
  \draw[-Latex, grigio] (0,-1.3) -- (0,1.35) node[above, black] {$y$};
  \node[font=\scriptsize, below] at (0.5,0) {$\frac12$};
  \node[font=\scriptsize, below right] at (2,0) {$2$};
  \node[font=\scriptsize, blu] at (0.5,0.27) {$C_{1}$};
  \node[font=\scriptsize, blu] at (1.85,0.85) {$C_{2}$};
\end{tikzpicture}
\caption{I domini delle funzioni \eqref{eq:domini-esempi}. A sinistra,
$(1-x^{2})(1-y^{2})\ge0$: il quadrato centrale e le quattro regioni d'angolo. A destra,
la lunula compresa tra la circonferenza $C_{1}$ di centro $(\frac12,0)$ e raggio
$\frac12$ e la circonferenza $C_{2}$ di centro $(1,0)$ e raggio $1$.}
\label{fig:domini}
\end{figure}

\begin{notastorica}[Le curve di livello: da Halley alle carte del tempo]
La rappresentazione di una grandezza che varia sulla superficie terrestre mediante linee
che uniscono i punti in cui essa ha lo stesso valore risale all'astronomo Edmond Halley,
che nel 1701 pubblicò una carta dell'Atlantico con le linee di uguale declinazione
magnetica, preziose per la navigazione con la bussola. Nel 1752 il geografo francese
Philippe Buache pubblicò una carta del canale della Manica con le linee di uguale
profondità, e alla
fine del Settecento le curve di livello del terreno entrarono nella cartografia
topografica. Nell'Ottocento l'idea passò alla meteorologia: le carte sinottiche con le
isobare permisero di riconoscere i cicloni e gli anticicloni (questo termine fu coniato
da Francis Galton nel 1863). Oggi la stessa idea si ritrova nelle carte delle
prestazioni dei velivoli e dei motori, e nelle visualizzazioni dei risultati delle
simulazioni fluidodinamiche (CFD), dove le linee di uguale pressione o di uguale velocità
si chiamano \emph{isolinee}.
\end{notastorica}

%% ----------------------------------------------------------------------------
%%  §7.2  Limiti e continuità
%% ----------------------------------------------------------------------------
\section{Limiti e continuità}
\label{sec:limiti-piano}

Per le funzioni di una variabile la vicinanza tra due numeri è misurata dal valore
assoluto della differenza; nel piano la misuriamo con la distanza
\eqref{eq:distanza}. Il ruolo degli intervalli $(x_{0}-\delta,x_{0}+\delta)$ è preso dai
cerchi.

\begin{definizione}[Intorno circolare, punto di accumulazione]\label{def:intorno-circolare}
\index{intorno circolare}\index{punto!di accumulazione}
Si chiama \emph{intorno circolare} di raggio $\delta>0$ del punto $P_{0}=(x_{0},y_{0})$ il
cerchio aperto di centro $P_{0}$ e raggio $\delta$:
\begin{equation}\label{eq:intorno-circolare}
\begin{split}
  I_{\delta}(P_{0})&=\set{(x,y)\in\R^{2} : (x-x_{0})^{2}+(y-y_{0})^{2}<\delta^{2}}\\
  &=\set{P\in\R^{2} : d(P,P_{0})<\delta}.
\end{split}
\end{equation}
Sia $D\subseteq\R^{2}$. Si dice che $P_{0}\in\R^{2}$ è un \emph{punto di accumulazione}
per $D$ se in ogni intorno circolare di $P_{0}$ cade almeno un punto di $D$ distinto da
$P_{0}$ (\cref{fig:intorni-piano}).
\end{definizione}

\begin{figure}[htb]
\centering
\begin{tikzpicture}[scale=0.9]
  \draw[-Latex, grigio] (-0.3,0) -- (4.2,0) node[right, black] {$x$};
  \draw[-Latex, grigio] (0,-0.3) -- (0,3.2) node[above, black] {$y$};
  \fill[blu!8] (2.2,1.5) circle (1.1);
  \draw[blu, thick, dashed] (2.2,1.5) circle (1.1);
  \fill (2.2,1.5) circle (1.3pt);
  \draw[tratteggio] (0,1.5) node[left, black, font=\scriptsize] {$y_{0}$} -- (2.2,1.5)
    -- (2.2,0) node[below, black, font=\scriptsize] {$x_{0}$};
  \draw[rame] (2.2,1.5) -- node[above, font=\scriptsize] {$\delta$} ++(40:1.1);
  \node[blu, font=\scriptsize] at (3.55,2.55) {$I_{\delta}(P_{0})$};
\end{tikzpicture}\hspace{1.2cm}
\begin{tikzpicture}[scale=0.9]
  \draw[-Latex, grigio] (-0.3,0) -- (4.2,0) node[right, black] {$x$};
  \draw[-Latex, grigio] (0,-0.3) -- (0,3.2) node[above, black] {$y$};
  \fill[pattern=north east lines, pattern color=blu]
    (1.8,1.2) .. controls (2.6,0.9) and (3.6,1.0) .. (3.8,1.9)
    .. controls (4.0,2.8) and (2.6,3.1) .. (2.1,2.6)
    .. controls (1.7,2.2) and (1.6,1.6) .. (1.8,1.2) -- cycle;
  \draw[blu, thick]
    (1.8,1.2) .. controls (2.6,0.9) and (3.6,1.0) .. (3.8,1.9)
    .. controls (4.0,2.8) and (2.6,3.1) .. (2.1,2.6)
    .. controls (1.7,2.2) and (1.6,1.6) .. (1.8,1.2) -- cycle;
  \foreach \r in {0.3,0.6,0.95} \draw[rame, dashed] (1.8,1.2) circle (\r);
  \fill (1.8,1.2) circle (1.3pt);
  \node[font=\scriptsize, fill=white, inner sep=1pt, anchor=north east] at (1.75,1.1)
    {$P_{0}$};
  \node[blu, font=\scriptsize] at (3.1,2.1) {$D$};
\end{tikzpicture}
\caption{A sinistra, l'intorno circolare $I_{\delta}(P_{0})$: il cerchio aperto di centro
$P_{0}=(x_{0},y_{0})$ e raggio $\delta$. A destra, un punto $P_{0}$ di accumulazione per
l'insieme $D$: ogni intorno circolare di $P_{0}$, per quanto piccolo, contiene punti di
$D$ diversi da $P_{0}$ (qui $P_{0}$ sta sul bordo di $D$).}
\label{fig:intorni-piano}
\end{figure}

\begin{definizione}[Limite di una funzione di due variabili]\label{def:limite-piano}
\index{limite!di una funzione di due variabili}
Sia $f:D\to\R$ e sia $P_{0}=(x_{0},y_{0})$ un punto di accumulazione per $D$. Si dice
che $f(x,y)$ \emph{converge al numero reale $\ell$} per $(x,y)$ che tende a
$(x_{0},y_{0})$, e si scrive
\begin{equation}\label{eq:limite-piano}
  \lim_{(x,y)\to(x_{0},y_{0})}f(x,y)=\ell,
\end{equation}
se per ogni $\varepsilon>0$ esiste $\delta>0$ tale che $\abs{f(x,y)-\ell}<\varepsilon$ per ogni
$(x,y)\in I_{\delta}(P_{0})\cap D$ diverso da $P_{0}$. In simboli:
\begin{equation}\label{eq:limite-piano-simboli}
\begin{gathered}
  \perogni\varepsilon>0\ \ \esiste\delta>0 :\quad \perogni(x,y)\in D,\\
  0<\sqrt{(x-x_{0})^{2}+(y-y_{0})^{2}}<\delta \ \implica\ \abs{f(x,y)-\ell}<\varepsilon .
\end{gathered}
\end{equation}
Analogamente si definiscono i limiti $\pm\infty$.
\end{definizione}

\begin{osservazione}[Limiti e successioni]\label{oss:limite-piano-successioni}
Una successione di punti $P_{n}=(x_{n},y_{n})$ converge a $P_{0}=(x_{0},y_{0})$ se
$d(P_{n},P_{0})\to0$; per le \eqref{eq:distanza-componenti}, ciò equivale a $x_{n}\to x_{0}$ e
$y_{n}\to y_{0}$. Con la stessa dimostrazione del \cref{teo:eps-delta}, nella quale basta
sostituire il valore assoluto con la distanza, si prova che la
\eqref{eq:limite-piano} equivale alla condizione: \emph{per ogni successione
$P_{n}\to P_{0}$ di punti di $D$ diversi da $P_{0}$, risulta $f(P_{n})\to\ell$}. Ne segue,
come per le funzioni di una variabile, che il limite, se esiste, è unico.
\end{osservazione}

\begin{esempio}\label{es:limite-x4-y4}
Proviamo che
\begin{equation}\label{eq:limite-x4-y4}
  \lim_{(x,y)\to(0,0)}\frac{x^{4}+y^{4}}{x^{2}+y^{2}}=0 .
\end{equation}
Poiché $x^{4}+y^{4}\le x^{4}+y^{4}+2x^{2}y^{2}=(x^{2}+y^{2})^{2}$, per $(x,y)\neq(0,0)$ si ha
\[
  0\le\frac{x^{4}+y^{4}}{x^{2}+y^{2}}\le\frac{(x^{2}+y^{2})^{2}}{x^{2}+y^{2}}=x^{2}+y^{2} .
\]
Fissato $\varepsilon>0$ e posto $\delta=\sqrt\varepsilon$, se $0<\sqrt{x^{2}+y^{2}}<\delta$ allora
$x^{2}+y^{2}<\varepsilon$, e quindi $\Bigl|\frac{x^{4}+y^{4}}{x^{2}+y^{2}}-0\Bigr|<\varepsilon$.
\end{esempio}

\begin{definizione}[Continuità]\label{def:continua-piano}\index{funzione!continua!in più variabili}
Sia $f:D\to\R$ e sia $P_{0}\in D$ un punto di accumulazione per $D$. Si dice che $f$ è
\emph{continua} in $P_{0}$ se
\begin{equation}\label{eq:continua-piano}
  \lim_{(x,y)\to(x_{0},y_{0})}f(x,y)=f(x_{0},y_{0}).
\end{equation}
Si dice che $f$ è continua in $D$ se è continua in ogni punto di $D$.
\end{definizione}

\begin{esempio}[Il prodotto $xy$ è continuo nell'origine]\label{es:xy-continua}
Verifichiamo che $f(x,y)=xy$ è continua nell'origine, cioè che $xy\to0$ per
$(x,y)\to(0,0)$. A tale scopo osserviamo che
\begin{equation}\label{eq:xy-disuguaglianza}
  \abs{xy}\le\frac12(x^{2}+y^{2}) \qquad\perogni(x,y)\in\R^{2}:
\end{equation}
infatti la disuguaglianza equivale a $x^{2}+y^{2}-2\abs x\abs y\ge0$, cioè a
$(\abs x-\abs y)^{2}\ge0$. Per la \eqref{eq:xy-disuguaglianza}, risulta $\abs{xy}<\varepsilon$
ogni volta che $x^{2}+y^{2}<2\varepsilon$: posto $\delta=\sqrt{2\varepsilon}$, da
$\sqrt{x^{2}+y^{2}}<\delta$ segue $\abs{xy-0}<\varepsilon$.
\end{esempio}

Come accade per le funzioni di una variabile, anche le funzioni di due variabili
ottenute componendo potenze, esponenziali, logaritmi e funzioni trigonometriche sono
continue nel loro insieme di definizione. Il motivo è il teorema seguente.

\begin{teorema}[Funzioni continue di due variabili]\label{teo:continue-piano}
\begin{enumerate}[label=\textup{(\roman*)}]
  \item Le funzioni $(x,y)\mapsto x$ e $(x,y)\mapsto y$ sono continue in $\R^{2}$.
  \item Somma, differenza e prodotto di funzioni continue in $P_{0}$ sono continui in
  $P_{0}$; lo è anche il quoziente, se il denominatore non si annulla in $P_{0}$.
  \item Se $f$ è continua in $P_{0}$ e $g$ è una funzione di una variabile continua nel
  punto $f(P_{0})$, allora la funzione composta $g\bigl(f(x,y)\bigr)$ è continua in $P_{0}$.
\end{enumerate}
\end{teorema}

\begin{schema}
\Ipotesi{(i) nessuna; \ (ii) $f$, $g$ continue in $P_{0}$; \ (iii) $f$ continua in
$P_{0}$, $g$ continua in $f(P_{0})$.}
\Tesi{continuità in $P_{0}$ delle funzioni coordinate, di $f\pm g$, $fg$, $f/g$, $g\circ f$.}
\end{schema}

\begin{strategia}
Grazie all'\cref{oss:limite-piano-successioni}, la continuità in $P_{0}$ equivale a
$f(P_{n})\to f(P_{0})$ per ogni successione $P_{n}\to P_{0}$. In questo modo le
proprietà delle funzioni di due variabili si riducono a quelle, già note, delle
successioni di numeri reali e delle funzioni continue di una variabile.
\end{strategia}

\begin{dimostrazione}
(i) Per la \eqref{eq:distanza-componenti}, $\abs{x-x_{0}}\le d(P,P_{0})$: se la distanza è
minore di $\delta=\varepsilon$, anche $\abs{x-x_{0}}<\varepsilon$. Lo stesso per $y$.
(ii) Sia $P_{n}\to P_{0}$, con $P_{n}\in D$. Per la continuità, $f(P_{n})\to f(P_{0})$ e
$g(P_{n})\to g(P_{0})$; per il teorema sulle operazioni con i limiti di successioni
$f(P_{n})\pm g(P_{n})\to f(P_{0})\pm g(P_{0})$, e analogamente per il prodotto e il
quoziente. Poiché la successione è arbitraria, si conclude con
l'\cref{oss:limite-piano-successioni}.
(iii) Se $P_{n}\to P_{0}$, allora $t_{n}=f(P_{n})\to f(P_{0})$; per la continuità di $g$
nel punto $f(P_{0})$ (\cref{prop:caratt-continuita}), $g(t_{n})\to g\bigl(f(P_{0})\bigr)$.
\end{dimostrazione}

Dai punti (i) e (ii) segue che ogni polinomio in $x$ e $y$ è continuo in $\R^{2}$; dal
punto (iii), che funzioni come $\cos(x^{2}+y^{2})$, $\sen(xy)$, $e^{-(x^{2}+y^{2})}$ sono
continue in $\R^{2}$, e che $\log(x^{2}+y^{2})$ è continua in $\R^{2}\setminus\{(0,0)\}$.

\subsection{Come si dimostra che un limite non esiste}

Nel piano ci si può avvicinare a un punto lungo infinite direzioni e lungo infinite
curve. Se il limite esiste, il risultato non può dipendere dal percorso.

\begin{proposizione}[Restrizione a una curva]\label{prop:restrizione-curva}
Sia $\lim_{(x,y)\to P_{0}}f(x,y)=\ell$, e siano $x(t)$, $y(t)$ funzioni definite per
$0<t<r$ tali che il punto $\gamma(t)=\bigl(x(t),y(t)\bigr)$ appartenga a $D$, sia diverso
da $P_{0}$ e tenda a $P_{0}$ per $t\to0^{+}$. Allora $f\bigl(x(t),y(t)\bigr)\to\ell$ per
$t\to0^{+}$.
\end{proposizione}

\begin{schema}
\Ipotesi{$f(x,y)\to\ell$ per $(x,y)\to P_{0}$; $\gamma(t)\in D\setminus\{P_{0}\}$,
$\gamma(t)\to P_{0}$ per $t\to0^{+}$.}
\Tesi{$f(\gamma(t))\to\ell$ per $t\to0^{+}$.}
\end{schema}

\begin{strategia}
Se $f$ è vicina a $\ell$ in tutti i punti di un intorno di $P_{0}$, lo è in particolare nei
punti della curva che cadono in quell'intorno; e la curva, tendendo a $P_{0}$, finisce per
restarci.
\end{strategia}

\begin{dimostrazione}
Fissato $\varepsilon>0$, sia $\delta$ come nella \eqref{eq:limite-piano-simboli}. Poiché
$\gamma(t)\to P_{0}$, esiste $\tau>0$ tale che $0<d\bigl(\gamma(t),P_{0}\bigr)<\delta$ per
$0<t<\tau$; per tali $t$ risulta $\abs{f(\gamma(t))-\ell}<\varepsilon$.
\end{dimostrazione}

In forma contronominale: \emph{se lungo due curve che tendono a $P_{0}$ la funzione ha
limiti diversi, o se lungo una curva non ha limite, il limite per $(x,y)\to P_{0}$ non
esiste}. È il corrispondente del \cref{cor:non-esistenza}.

\begin{esempio}[Una funzione con limiti diversi lungo le rette]\label{es:xy-su-x2y2}
Sia $f(x,y)=\frac{xy}{x^{2}+y^{2}}$ per $(x,y)\neq(0,0)$. Lungo la retta $y=mx$, con
$x\neq0$,
\[
  f(x,mx)=\frac{mx^{2}}{x^{2}+m^{2}x^{2}}=\frac{m}{1+m^{2}},
\]
che non dipende da $x$: la funzione è costante su ogni retta per l'origine, ma con un
valore che cambia con la pendenza $m$ (vale $\frac12$ sulla retta $y=x$, $-\frac12$ sulla retta
$y=-x$, $0$ sugli assi). Quindi il limite per $(x,y)\to(0,0)$ non esiste
(\cref{fig:discontinua}). In coordinate polari, $x=\rho\cos\vartheta$,
$y=\rho\sen\vartheta$, si ha $f=\cos\vartheta\sen\vartheta=\frac12\sen2\vartheta$: il valore
dipende soltanto dall'angolo.
\end{esempio}

\begin{figure}[htb]
\centering
\includegraphics[width=0.95\textwidth]{discontinua}
\caption{La funzione $f(x,y)=\frac{xy}{x^{2}+y^{2}}$ è costante su ogni semiretta uscente
dall'origine (a sinistra, il colore indica il valore). A destra, il valore
$\frac12\sen2\vartheta$ sulla semiretta di angolo $\vartheta$. In ogni intorno dell'origine,
come quello tratteggiato, la funzione assume tutti i valori di $[-\frac12,\frac12]$: non ha
limite per $(x,y)\to(0,0)$ (\texttt{python/fig\_discontinua.py}).}
\label{fig:discontinua}
\end{figure}

\begin{attenzione}[Le rette non bastano]
Si potrebbe pensare che, se la funzione ha lo stesso limite lungo \emph{tutte} le rette
per $P_{0}$, il limite esista. Non è così. Sia $g(x,y)=\frac{x^{2}y}{x^{4}+y^{2}}$ per
$(x,y)\neq(0,0)$. Lungo l'asse $x=0$ la funzione è nulla; lungo la retta $y=mx$, con
$m\neq0$,
\[
  g(x,mx)=\frac{mx^{3}}{x^{4}+m^{2}x^{2}}=\frac{mx}{x^{2}+m^{2}}\ \to\ 0 \qquad (x\to0).
\]
Ma lungo la parabola $y=x^{2}$ si ha $g(x,x^{2})=\frac{x^{4}}{2x^{4}}=\frac12$ per ogni
$x\neq0$: il limite non esiste. Il programma seguente mostra i valori di $f$ e di $g$
lungo alcuni percorsi.
\end{attenzione}

\begin{python}[Valori lungo percorsi diversi verso l'origine]
f = lambda x, y: x * y / (x**2 + y**2)
g = lambda x, y: x**2 * y / (x**4 + y**2)

print("   t      f(t,0)  f(t,t)  f(t,2t)"
      "    g(t,t)  g(t,3t)  g(t,t^2)")
for t in (0.1, 0.01, 0.001, 0.0001):
    print(f"{t:7.4f}  {f(t, 0):6.3f}  {f(t, t):6.3f}"
          f"  {f(t, 2*t):6.3f}   {g(t, t):7.4f}"
          f"  {g(t, 3*t):7.4f}  {g(t, t**2):7.4f}")
\end{python}
\begin{risultato}
   t      f(t,0)  f(t,t)  f(t,2t)    g(t,t)  g(t,3t)  g(t,t^2)
 0.1000   0.000   0.500   0.400    0.0990   0.0333   0.5000
 0.0100   0.000   0.500   0.400    0.0100   0.0033   0.5000
 0.0010   0.000   0.500   0.400    0.0010   0.0003   0.5000
 0.0001   0.000   0.500   0.400    0.0001   0.0000   0.5000
\end{risultato}

Per dimostrare che un limite \emph{esiste} serve dunque un argomento che controlli tutti
i punti di un intorno contemporaneamente, come nell'\cref{es:limite-x4-y4}. Uno strumento
comodo è il passaggio alle coordinate polari.

\begin{proposizione}[Criterio delle coordinate polari]\label{prop:coordinate-polari}
Sia $f$ definita in un intorno di $P_{0}=(x_{0},y_{0})$, eccetto al più in $P_{0}$, e sia
$\ell\in\R$. Se esiste una funzione $\omega(\rho)$, definita per $\rho>0$ piccolo, tale che
$\omega(\rho)\to0$ per $\rho\to0^{+}$ e
\begin{equation}\label{eq:coordinate-polari}
  \bigl|f(x_{0}+\rho\cos\vartheta,\,y_{0}+\rho\sen\vartheta)-\ell\bigr|\le\omega(\rho)
  \qquad\text{per ogni }\vartheta\in[0,2\pi),
\end{equation}
allora $\lim_{(x,y)\to P_{0}}f(x,y)=\ell$.
\end{proposizione}

\begin{schema}
\Ipotesi{una maggiorazione di $\abs{f-\ell}$ che dipende soltanto dalla distanza $\rho$
da $P_{0}$, \emph{non} dall'angolo $\vartheta$, e tende a $0$.}
\Tesi{$f(x,y)\to\ell$ per $(x,y)\to P_{0}$.}
\end{schema}

\begin{strategia}
Ogni punto $P\neq P_{0}$ si scrive come $P_{0}+\rho(\cos\vartheta,\sen\vartheta)$ con
$\rho=d(P,P_{0})$. La maggiorazione, valendo per ogni angolo, controlla allora $f$ su
tutta la circonferenza di raggio $\rho$, e quindi su tutto l'intorno circolare.
\end{strategia}

\begin{dimostrazione}
Fissato $\varepsilon>0$, sia $\delta>0$ tale che $\omega(\rho)<\varepsilon$ per $0<\rho<\delta$.
Se $0<d(P,P_{0})<\delta$, posto $\rho=d(P,P_{0})$ esiste $\vartheta$ con
$P=(x_{0}+\rho\cos\vartheta,\,y_{0}+\rho\sen\vartheta)$, e per la \eqref{eq:coordinate-polari}
$\abs{f(P)-\ell}\le\omega(\rho)<\varepsilon$.
\end{dimostrazione}

Ad esempio, per $f(x,y)=\frac{x^{2}y}{x^{2}+y^{2}}$ si ha
$\abs{f(\rho\cos\vartheta,\rho\sen\vartheta)}=\rho\abs{\cos^{2}\vartheta\sen\vartheta}\le\rho$, e
quindi $f(x,y)\to0$ per $(x,y)\to(0,0)$. È essenziale che la maggiorazione sia
indipendente da $\vartheta$: per la funzione $g$ dell'avvertenza precedente si ha
$g=\frac{\rho\cos^{2}\vartheta\sen\vartheta}{\rho^{2}\cos^{4}\vartheta+\sen^{2}\vartheta}$, che
per ogni angolo \emph{fissato} tende a $0$, ma non uniformemente rispetto a $\vartheta$.

\begin{osservazione}[Il teorema di Weierstrass nel piano]\label{oss:weierstrass-piano}
Si dimostra (nei corsi di Analisi Matematica di secondo anno) che una funzione continua
su un insieme \emph{chiuso e limitato} del piano, come un rettangolo
$[a,b]\times[c,d]$ o un cerchio insieme alla sua circonferenza, assume massimo e minimo
assoluti: è il teorema di Weierstrass in più variabili, che useremo come il
\cref{teo:weierstrass} nella ricerca degli estremi.
\end{osservazione}

\begin{notastorica}[Dalla distanza agli spazi metrici: Fréchet e Hausdorff]
In questo paragrafo abbiamo ricopiato le definizioni di limite e di continuità date per
le funzioni di una variabile, sostituendo il valore assoluto $\abs{x-x_{0}}$ con la
distanza $d(P,P_{0})$. Che l'operazione riesca non è un caso: quasi tutte le
dimostrazioni usano soltanto tre proprietà della distanza, cioè che $d(P,Q)\ge0$ con
l'uguaglianza solo per $P=Q$, che $d(P,Q)=d(Q,P)$, e la disuguaglianza triangolare.
L'idea di fare analisi in «spazi» diversi dalla retta risale a Bernhard Riemann, che
nella lezione di abilitazione del 1854 \emph{Sulle ipotesi che stanno alla base della
geometria} parlò senza esitare di varietà con infinite dimensioni, come l'insieme delle
funzioni definite in un dominio; più tardi Georg Cantor introdusse i punti di
accumulazione e gli insiemi aperti e chiusi. Il passo decisivo fu di
Maurice Fréchet, nella tesi \emph{Sur quelques points du calcul fonctionnel} (1906): in
un insieme qualsiasi, purché vi sia definito tra due elementi uno «scarto»
(\emph{écart}) con quelle proprietà, si possono definire limiti, continuità e
compattezza (si veda anche la nota storica del \S\ref{sec:lineari}). Il nome di
\emph{spazio metrico} si deve a Felix Hausdorff, che nei \emph{Grundzüge der
Mengenlehre} (1914) osservò inoltre che molte nozioni dipendono soltanto dalla famiglia
degli intorni, e non dalla distanza che li genera: è la nascita degli spazi topologici.
La lezione, per noi, è concreta: ciò che si dimostra con la sola disuguaglianza
triangolare vale nel piano, in $\R^{n}$ (\cref{sec:n-variabili}) e negli spazi di
funzioni con cui lavorano i metodi numerici.
\end{notastorica}

%% ----------------------------------------------------------------------------
%%  §7.3  Derivate parziali. Gradiente
%% ----------------------------------------------------------------------------
\section{Derivate parziali. Gradiente}
\label{sec:derivate-parziali}

Sia $f$ una funzione di due variabili definita in un intorno circolare $I_{\delta}$ di un
punto $(x,y)$. In particolare $f$ è definita, oltre che in $(x,y)$, anche nei punti
$(x+h,y)$ con $\abs h<\delta$, che stanno sulla retta orizzontale per $(x,y)$, e nei punti
$(x,y+k)$ con $\abs k<\delta$, sulla retta verticale. Fermando una delle due variabili si
ottiene una funzione di una sola variabile, che possiamo derivare.

\begin{definizione}[Derivate parziali]\label{def:derivate-parziali}\index{derivata!parziale}
La \emph{derivata parziale di $f$ rispetto a $x$} nel punto $(x,y)$ è il limite
\begin{equation}\label{eq:derivata-parziale-x}
  \lim_{h\to0}\frac{f(x+h,y)-f(x,y)}{h},
\end{equation}
se esiste ed è finito; si indica con uno dei simboli
\begin{equation}\label{eq:notazioni-parziale-x}
  f_{x}, \qquad f_{x}(x,y), \qquad \frac{\partial f}{\partial x}, \qquad D_{x}f .
\end{equation}
Analogamente, la \emph{derivata parziale di $f$ rispetto a $y$} nel punto $(x,y)$ è il
limite
\begin{equation}\label{eq:derivata-parziale-y}
  \lim_{k\to0}\frac{f(x,y+k)-f(x,y)}{k},
\end{equation}
se esiste ed è finito, e si indica con $f_{y}$, $f_{y}(x,y)$, $\frac{\partial f}{\partial y}$,
$D_{y}f$.
\end{definizione}

Si noti che la \eqref{eq:derivata-parziale-x} è il limite del rapporto incrementale di una
funzione di \emph{una} variabile: posto $g(t)=f(t,y)$, con $y$ fissato che gioca il ruolo
di parametro, è $f_{x}(x,y)=g'(x)$. Per questo motivo \emph{le derivate parziali si
calcolano con le stesse regole di derivazione delle funzioni di una variabile,
considerando l'altra variabile costante}. Geometricamente, $f_{x}(x_{0},y_{0})$ è la
pendenza della curva che si ottiene sezionando il grafico con il piano verticale
$y=y_{0}$, e $f_{y}(x_{0},y_{0})$ quella della sezione con il piano $x=x_{0}$
(\cref{fig:derivate-parziali}).

\begin{figure}[htb]
\centering
\includegraphics[width=0.62\textwidth]{derivate_parziali}
\caption{Il significato geometrico delle derivate parziali per
$f(x,y)=3-\frac{x^{2}}2-\frac{y^{2}}4$ nel punto $(x_{0},y_{0})=(1,1)$. La sezione con il
piano $y=y_{0}$ (rame) ha nel punto $(x_{0},y_{0},f(x_{0},y_{0}))$ tangente di pendenza
$f_{x}(1,1)=-1$; la sezione con il piano $x=x_{0}$ (verde) ha tangente di pendenza
$f_{y}(1,1)=-\frac12$ (\texttt{python/fig\_derivate\_parziali.py}).}
\label{fig:derivate-parziali}
\end{figure}

\begin{esempio}[Calcolo di derivate parziali]\label{es:derivate-parziali}
Per la funzione
\begin{equation}\label{eq:parziali-es1}
  f(x,y)=x^{2}+xy^{2}, \qquad (x,y)\in\R^{2},
\end{equation}
nel derivare rispetto a $x$ si considera $y$ costante, come se fosse, ad esempio,
$y^{2}=9$ e si dovesse derivare la funzione $x\mapsto x^{2}+9x$; si ottiene
$f_{x}(x,y)=2x+y^{2}$. Analogamente, derivando rispetto a $y$ con $x$ costante,
$f_{y}(x,y)=2xy$. La funzione $f(x,y)=\sen(xy)$, il cui grafico è nella
\cref{fig:superfici}, ha derivate parziali $f_{x}=y\cos(xy)$ e $f_{y}=x\cos(xy)$ in ogni
punto di $\R^{2}$ (regola della catena). La funzione $f(x,y)=e^{x}$ è costante rispetto a
$y$: le sue derivate parziali sono $f_{x}=e^{x}$ e $f_{y}=0$.
\end{esempio}

\begin{python}[Derivate parziali con sympy e controllo numerico]
import sympy as sp

x, y = sp.symbols("x y", real=True)
for f in (x**2 + x*y**2, sp.sin(x*y), sp.exp(x), x + 2*y):
    gx, gy = sp.diff(f, x), sp.diff(f, y)
    print(f"f = {str(f):13}  f_x = {str(gx):10}  f_y = {gy}")

# controllo numerico con i rapporti incrementali, nel punto (1, 2)
F = sp.lambdify((x, y), x**2 + x*y**2)
h = 1e-6
print("f_x(1,2) ~", round((F(1 + h, 2) - F(1, 2)) / h, 4),
      "  f_y(1,2) ~", round((F(1, 2 + h) - F(1, 2)) / h, 4))
\end{python}
\begin{risultato}
f = x**2 + x*y**2  f_x = 2*x + y**2  f_y = 2*x*y
f = sin(x*y)       f_x = y*cos(x*y)  f_y = x*cos(x*y)
f = exp(x)         f_x = exp(x)      f_y = 0
f = x + 2*y        f_x = 1           f_y = 2
f_x(1,2) ~ 6.0   f_y(1,2) ~ 4.0
\end{risultato}

\begin{esempio}[Sensibilità della resistenza all'allungamento alare]\label{es:sensibilita-AR}
Con la polare parabolica, il coefficiente di resistenza dipende dal coefficiente di
portanza e dall'allungamento alare:
\begin{equation}\label{eq:CD-CL-AR}
  C_{D}(C_{L},AR)=C_{D0}+\frac{C_{L}^{2}}{\pi e\,AR} .
\end{equation}
Le derivate parziali misurano la \emph{sensibilità} di $C_{D}$ a ciascuna delle due
variabili, a parità dell'altra:
\[
  \frac{\partial C_{D}}{\partial C_{L}}=\frac{2C_{L}}{\pi e\,AR}, \qquad
  \frac{\partial C_{D}}{\partial AR}=-\frac{C_{L}^{2}}{\pi e\,AR^{2}} .
\]
La seconda è negativa: aumentando l'allungamento la resistenza indotta diminuisce, ma
sempre meno, perché il modulo della derivata decresce come $\frac1{AR^{2}}$; ed è tanto
maggiore in modulo quanto più alto è $C_{L}$, cioè nel volo lento, in salita e in
virata. È il motivo per cui alianti e velivoli da alta quota hanno ali molto allungate.
\end{esempio}

\begin{python}[Sensibilità di $C_{D}$ a $C_{L}$ e ad $AR$]
import math

# C_D(C_L, AR) = C_D0 + C_L^2 / (pi e AR): sensibilità parziali
CD0, e = 0.025, 0.8
CD = lambda CL, AR: CD0 + CL**2 / (math.pi * e * AR)
dCD_dCL = lambda CL, AR: 2 * CL / (math.pi * e * AR)
dCD_dAR = lambda CL, AR: -CL**2 / (math.pi * e * AR**2)
print("  C_L  AR     C_D     dC_D/dC_L   dC_D/dAR")
for CL, AR in ((0.5, 8), (1.0, 8), (1.0, 12)):
    print(f"  {CL:.1f}  {AR:2d}   {CD(CL, AR):.5f}   "
          f"{dCD_dCL(CL, AR):.5f}    {dCD_dAR(CL, AR):.6f}")
\end{python}
\begin{risultato}
  C_L  AR     C_D     dC_D/dC_L   dC_D/dAR
  0.5   8   0.03743   0.04974    -0.001554
  1.0   8   0.07474   0.09947    -0.006217
  1.0  12   0.05816   0.06631    -0.002763
\end{risultato}

\subsection{Il gradiente e le derivate direzionali}

\begin{definizione}[Gradiente]\label{def:gradiente}\index{gradiente}
Se $f$ ammette le derivate parziali $f_{x}$, $f_{y}$ nel punto $(x,y)$, si chiama
\emph{gradiente} di $f$ in $(x,y)$ il vettore di $\R^{2}$ che ha per componenti le
derivate parziali:
\begin{equation}\label{eq:gradiente}
  \nabla f(x,y)=\operatorname{grad}f(x,y)=\bigl(f_{x}(x,y),\,f_{y}(x,y)\bigr).
\end{equation}
\end{definizione}

Il simbolo $\nabla$ si legge «nabla» (dal nome di un'arpa fenicia, di forma triangolare).
Le derivate parziali descrivono come varia $f$ quando ci si sposta parallelamente agli
assi; ci si può però spostare in qualunque direzione. Sia $v=(\cos\vartheta,\sen\vartheta)$
un \emph{versore}, cioè un vettore di lunghezza $1$; la \emph{derivata direzionale} di $f$
nel punto $P=(x,y)$ secondo la direzione $v$ è il limite
\begin{equation}\label{eq:derivata-direzionale}
  \frac{\partial f}{\partial v}(x,y)=\lim_{t\to0}\frac{f(x+t\cos\vartheta,\,y+t\sen\vartheta)-f(x,y)}{t},
\end{equation}
se esiste finito. Per $\vartheta=0$ si ritrova $f_{x}$, per $\vartheta=\frac\pi2$ si ritrova
$f_{y}$. Il gradiente ha una proprietà notevole: \emph{se non è nullo, indica la direzione
di massima pendenza del grafico}. Lo verifichiamo qui per le funzioni lineari; il caso
generale, per le funzioni differenziabili, è trattato nel \cref{sec:differenziabilita}.

\begin{proposizione}[Gradiente di una funzione lineare]\label{prop:gradiente-lineare}
Sia $f(x,y)=ax+by+c$, con $(a,b)\neq(0,0)$. Allora $\nabla f=(a,b)$ in ogni punto, e per
ogni versore $v=(\cos\vartheta,\sen\vartheta)$
\begin{equation}\label{eq:direzionale-lineare}
  \frac{\partial f}{\partial v}=a\cos\vartheta+b\sen\vartheta=\abs{\nabla f}\cos(\vartheta-\varphi),
\end{equation}
dove $\abs{\nabla f}=\sqrt{a^{2}+b^{2}}$ e $\varphi$ è l'angolo che il gradiente forma con
l'asse $x$. In particolare la derivata direzionale è massima, e uguale a $\abs{\nabla f}$,
nella direzione del gradiente; è minima nella direzione opposta; è nulla nelle direzioni
perpendicolari al gradiente, che sono quelle delle curve di livello.
\end{proposizione}

\begin{schema}
\Ipotesi{$f(x,y)=ax+by+c$, $(a,b)\neq(0,0)$; $v=(\cos\vartheta,\sen\vartheta)$.}
\Tesi{$\frac{\partial f}{\partial v}=\abs{\nabla f}\cos(\vartheta-\varphi)$: massima per
$v$ parallelo e concorde a $\nabla f$, nulla per $v$ perpendicolare.}
\end{schema}

\begin{strategia}
Per una funzione lineare il rapporto incrementale è costante, e la derivata direzionale
si calcola subito. Scrivendo il gradiente in forma polare,
$(a,b)=\abs{\nabla f}(\cos\varphi,\sen\varphi)$, la formula di sottrazione del coseno
trasforma $a\cos\vartheta+b\sen\vartheta$ in un solo coseno, di cui si conoscono massimo e
minimo.
\end{strategia}

\begin{dimostrazione}
Si ha $f_{x}=a$ e $f_{y}=b$. Inoltre
$f(x+t\cos\vartheta,y+t\sen\vartheta)-f(x,y)=t(a\cos\vartheta+b\sen\vartheta)$, quindi il rapporto
incrementale in \eqref{eq:derivata-direzionale} è costante e uguale a
$a\cos\vartheta+b\sen\vartheta$. Poiché $(a,b)\neq(0,0)$, posto $r=\sqrt{a^{2}+b^{2}}$ esiste
$\varphi$ con $a=r\cos\varphi$ e $b=r\sen\varphi$; per la formula di sottrazione
\eqref{eq:addizione-cos},
\[
  a\cos\vartheta+b\sen\vartheta=r(\cos\varphi\cos\vartheta+\sen\varphi\sen\vartheta)=r\cos(\vartheta-\varphi).
\]
Il massimo $r$ si ha per $\vartheta=\varphi$, il minimo $-r$ per $\vartheta=\varphi+\pi$, e il
valore $0$ per $\vartheta=\varphi\pm\frac\pi2$. Infine, le curve di livello
$ax+by+c=\text{costante}$ sono rette di direzione $(-b,a)$, perpendicolare a $(a,b)$.
\end{dimostrazione}

Per esempio, la funzione $f(x,y)=x+2y$ ha gradiente costante $(1,2)$: questo vettore
indica la direzione e il verso, nel piano di base, in cui conviene muoversi per ottenere
il massimo incremento di $f$ a parità di percorso. Il grafico di una funzione lineare è
un piano; per una funzione qualunque, nel \cref{sec:differenziabilita} vedremo che
vicino a ogni punto il grafico si confonde con un piano (il \emph{piano tangente}), e
che per questo la proprietà vale in generale.

\begin{idea}[La forza di gradiente di pressione]
Su una particella d'aria agisce, per effetto delle differenze di pressione, una forza
per unità di massa uguale a $-\frac1\rho\nabla p$: la forza è diretta come l'opposto del
gradiente, cioè dalle alte verso le basse pressioni, perpendicolarmente alle isobare, ed è
tanto più intensa quanto più le isobare sono fitte (\cref{fig:livello}, a sinistra). Se la
Terra non ruotasse, il vento soffierebbe proprio così; nella realtà la forza di Coriolis
lo devia, e alle quote alte il vento scorre quasi parallelo alle isobare (vento
geostrofico). La stessa forza governa lo strato limite sul profilo di un'ala: dove la
pressione cresce nel verso del moto (\emph{gradiente avverso}) la corrente rallenta e può
separarsi dalla superficie, provocando lo stallo.
\end{idea}

\begin{notastorica}[La nascita delle derivate parziali e del simbolo $\partial$]
Le derivate parziali nacquero insieme al calcolo, ma si affermarono negli anni Trenta e
Quaranta del Settecento, nei lavori di Eulero, di Clairaut e di d'Alembert sulla
meccanica dei corpi continui: nel 1747 d'Alembert scrisse l'equazione della corda
vibrante, la prima equazione alle derivate parziali della fisica matematica. Nessuno,
però, si fermò a darne una definizione esplicita: le si calcolava trattando come
costanti le altre variabili, esattamente come facciamo noi, e soltanto la revisione
critica dell'Ottocento rivelò che cosa esse non dicono
(\cref{sec:differenziabilita}). Per molto tempo si usò lo stesso simbolo $d$ delle
derivate ordinarie, specificando a parole la variabile. Il simbolo $\partial$ compare
già in una memoria di Condorcet del 1770; nella forma $\frac{\partial u}{\partial x}$ fu
usato da Adrien-Marie Legendre nel 1786, poi abbandonato, e reintrodotto in modo
definitivo da Carl Gustav Jacobi nel 1841. Il simbolo $\nabla$ e il nome «nabla» si
affermarono nell'Ottocento con il calcolo vettoriale di Hamilton, Tait, Maxwell e
Gibbs (il nome fu suggerito a Tait, intorno al 1870, dal biblista William Robertson
Smith).
\end{notastorica}

%% ----------------------------------------------------------------------------
%%  §7.4  Derivate successive. Teorema di Schwarz
%% ----------------------------------------------------------------------------
\section{Derivate successive. Teorema di Schwarz}
\label{sec:schwarz}

Sia $f$ una funzione di due variabili definita in un intorno circolare $I_{\delta}$ di un
punto di $\R^{2}$, e supponiamo che in tutti i punti di $I_{\delta}$ esistano le derivate
parziali $f_{x}(x,y)$ e $f_{y}(x,y)$. Queste sono a loro volta funzioni di due
variabili, e possono ammettere derivate parziali:
\begin{equation}\label{eq:derivate-seconde-def}
  \frac{\partial}{\partial x}f_{x}, \qquad \frac{\partial}{\partial y}f_{x}, \qquad
  \frac{\partial}{\partial x}f_{y}, \qquad \frac{\partial}{\partial y}f_{y} .
\end{equation}
Esse si chiamano \emph{derivate parziali seconde} di $f$ e si indicano rispettivamente con
i simboli
\begin{equation}\label{eq:derivate-seconde-notazioni}
\begin{gathered}
  f_{xx}, \qquad f_{xy}, \qquad f_{yx}, \qquad f_{yy},\\
  \text{oppure}\qquad
  \frac{\partial^{2}f}{\partial x^{2}}, \quad \frac{\partial^{2}f}{\partial x\,\partial y}, \quad
  \frac{\partial^{2}f}{\partial y\,\partial x}, \quad \frac{\partial^{2}f}{\partial y^{2}} .
\end{gathered}
\end{equation}
Le $f_{xy}$, $f_{yx}$ si dicono \emph{derivate seconde miste}, le $f_{xx}$, $f_{yy}$
\emph{derivate seconde pure}. Nella notazione con i pedici l'ordine di derivazione si
legge da sinistra a destra: $f_{xy}=(f_{x})_{y}$ si ottiene derivando prima rispetto a $x$
e poi rispetto a $y$.

\begin{attenzione}[Una questione di convenzioni]
Nella notazione $\frac{\partial^{2}f}{\partial x\,\partial y}$ l'ordine di derivazione non è
uniforme nella letteratura: molti testi, interpretando il simbolo come
$\frac{\partial}{\partial x}\bigl(\frac{\partial f}{\partial y}\bigr)$, lo usano per $f_{yx}$.
Nel seguito useremo soprattutto i pedici, il cui significato è univoco; vedremo comunque
che, sotto ipotesi molto generali, le due derivate miste coincidono e la questione
diventa irrilevante.
\end{attenzione}

\begin{esempio}[Le derivate seconde di $x^{y}$]\label{es:derivate-xy}
La funzione $f(x,y)=x^{y}$, definita nel semipiano $D=\set{(x,y)\in\R^{2} : x>0}$, per $y$
fissato è una potenza, mentre per $x$ fissato è un esponenziale, che ai fini della
derivazione rispetto a $y$ conviene scrivere nella forma $f(x,y)=e^{y\log x}$. Le derivate
parziali prime valgono
\begin{equation}\label{eq:xy-prime}
  f_{x}=y\,x^{y-1}, \qquad f_{y}=x^{y}\log x ;
\end{equation}
le derivate seconde pure sono $f_{xx}=y(y-1)x^{y-2}$ e $f_{yy}=x^{y}(\log x)^{2}$. Per
ottenere $f_{xy}$ deriviamo $f_{x}$ rispetto a $y$ con la regola del prodotto, e per
$f_{yx}$ deriviamo $f_{y}$ rispetto a $x$:
\begin{equation}\label{eq:xy-miste}
\begin{aligned}
  f_{xy}&=x^{y-1}+y\,x^{y-1}\log x=x^{y-1}(1+y\log x),\\
  f_{yx}&=y\,x^{y-1}\log x+x^{y}\cdot\frac1x=x^{y-1}(y\log x+1).
\end{aligned}
\end{equation}
\end{esempio}

Il lettore avrà notato che, con passaggi intermedi differenti, si è ottenuto lo stesso
risultato. Ciò non è casuale ed è formalizzato nell'importante teorema che segue.

\begin{teorema}[di Schwarz]\label{teo:schwarz}\index{teorema!di Schwarz}
Sia $f$ una funzione definita in un intorno circolare $I_{\delta}$ di $(x_{0},y_{0})$, nel
quale esistono le derivate $f_{x}$, $f_{y}$, $f_{xy}$, $f_{yx}$. Se le derivate miste
$f_{xy}$ e $f_{yx}$ sono continue in $(x_{0},y_{0})$, allora
\begin{equation}\label{eq:schwarz}
  f_{xy}(x_{0},y_{0})=f_{yx}(x_{0},y_{0}) .
\end{equation}
\end{teorema}

\begin{schema}
\Ipotesi{\begin{itemize}[nosep, leftmargin=1.2em]
  \item $f_{x}$, $f_{y}$, $f_{xy}$, $f_{yx}$ esistono in $I_{\delta}(x_{0},y_{0})$;
  \item $f_{xy}$ e $f_{yx}$ sono continue in $(x_{0},y_{0})$.
\end{itemize}}
\Tesi{$f_{xy}(x_{0},y_{0})=f_{yx}(x_{0},y_{0})$.}
\end{schema}

\begin{strategia}
Consideriamo il rettangolo di vertici opposti $(x_{0},y_{0})$ e $(x,y)$
(\cref{fig:schwarz}) e la combinazione dei valori di $f$ nei quattro vertici
\[
  \Delta=f(x,y)-f(x,y_{0})-f(x_{0},y)+f(x_{0},y_{0}),
\]
una specie di «rapporto incrementale doppio». La stessa quantità $\Delta$ si può leggere
in due modi: come incremento, tra $x_{0}$ e $x$, della funzione $F(s)=f(s,y)-f(s,y_{0})$,
oppure come incremento, tra $y_{0}$ e $y$, della funzione $G(t)=f(x,t)-f(x_{0},t)$. Applicando
due volte il teorema di Lagrange nel primo modo si trova
$\Delta=f_{xy}(x_{1},y_{1})(x-x_{0})(y-y_{0})$; nel secondo modo
$\Delta=f_{yx}(x_{2},y_{2})(x-x_{0})(y-y_{0})$. Quindi il valore di $f_{xy}$ in un punto del
rettangolo coincide con quello di $f_{yx}$ in un altro punto del rettangolo; facendo tendere il rettangolo al punto
$(x_{0},y_{0})$, la continuità delle derivate miste fa il resto.
\end{strategia}

\begin{dimostrazione}
Sia $(x,y)\in I_{\delta}$ con $x\neq x_{0}$ e $y\neq y_{0}$. Il rettangolo chiuso di vertici
opposti $(x_{0},y_{0})$ e $(x,y)$ è contenuto in $I_{\delta}$: infatti ogni suo punto
$(s,t)$ verifica $\abs{s-x_{0}}\le\abs{x-x_{0}}$ e $\abs{t-y_{0}}\le\abs{y-y_{0}}$, e quindi
dista da $(x_{0},y_{0})$ non più di $(x,y)$. Definiamo
\begin{equation}\label{eq:schwarz-FG}
\begin{aligned}
  F(s)&=f(s,y)-f(s,y_{0}) \qquad (y \text{ fissato}),\\
  G(t)&=f(x,t)-f(x_{0},t) \qquad (x \text{ fissato}).
\end{aligned}
\end{equation}
Risulta
\[
\begin{aligned}
  F(x)-F(x_{0})&=f(x,y)-f(x,y_{0})-\bigl[f(x_{0},y)-f(x_{0},y_{0})\bigr],\\
  G(y)-G(y_{0})&=f(x,y)-f(x_{0},y)-\bigl[f(x,y_{0})-f(x_{0},y_{0})\bigr],
\end{aligned}
\]
per cui
\begin{equation}\label{eq:schwarz-uguaglianza}
  F(x)-F(x_{0})=G(y)-G(y_{0}) .
\end{equation}
La funzione $F$ è derivabile nell'intervallo di estremi $x_{0}$, $x$, con
$F'(s)=f_{x}(s,y)-f_{x}(s,y_{0})$. Per il teorema di Lagrange (\cref{teo:lagrange}) esiste
$x_{1}$ nell'intervallo di estremi $x_{0}$, $x$ tale che
\[
  F(x)-F(x_{0})=F'(x_{1})(x-x_{0})=\bigl[f_{x}(x_{1},y)-f_{x}(x_{1},y_{0})\bigr](x-x_{0}).
\]
Applicando di nuovo il teorema di Lagrange alla funzione $t\mapsto f_{x}(x_{1},t)$, che
è derivabile con derivata $f_{xy}(x_{1},t)$, nell'intervallo di estremi $y_{0}$, $y$,
otteniamo un punto $y_{1}$ di tale intervallo per cui
\begin{equation}\label{eq:schwarz-F}
  F(x)-F(x_{0})=f_{xy}(x_{1},y_{1})(x-x_{0})(y-y_{0}) .
\end{equation}
Procedendo in modo analogo con la funzione $G$ nell'intervallo di estremi $y_{0}$, $y$,
troviamo $y_{2}$ e poi $x_{2}$ tali che
\begin{equation}\label{eq:schwarz-G}
\begin{split}
  G(y)-G(y_{0})&=\bigl[f_{y}(x,y_{2})-f_{y}(x_{0},y_{2})\bigr](y-y_{0})\\
  &=f_{yx}(x_{2},y_{2})(x-x_{0})(y-y_{0}).
\end{split}
\end{equation}
Confrontando le \eqref{eq:schwarz-uguaglianza}, \eqref{eq:schwarz-F}, \eqref{eq:schwarz-G}
e dividendo per $(x-x_{0})(y-y_{0})\neq0$ si ottiene
\begin{equation}\label{eq:schwarz-punti}
  f_{xy}(x_{1},y_{1})=f_{yx}(x_{2},y_{2}),
\end{equation}
dove i punti $(x_{1},y_{1})$ e $(x_{2},y_{2})$ appartengono al rettangolo. Scegliamo ora
$(x,y)=(x_{0}+\frac rn,\,y_{0}+\frac rn)$, con $r>0$ fissato abbastanza piccolo e
$n\in\N$, $n\ge1$. I punti corrispondenti $(x_{1},y_{1})$, $(x_{2},y_{2})$ distano da
$(x_{0},y_{0})$ meno di $\frac{\sqrt2\,r}{n}$, e quindi tendono a $(x_{0},y_{0})$ per
$n\to+\infty$. Per la continuità di $f_{xy}$ e $f_{yx}$ in $(x_{0},y_{0})$
(\cref{oss:limite-piano-successioni}), passando al limite nella \eqref{eq:schwarz-punti}
si ottiene la tesi.
\end{dimostrazione}

\begin{figure}[htb]
\centering
\begin{tikzpicture}[scale=1.0]
  \draw[-Latex, grigio] (-0.3,0) -- (6.2,0) node[right, black] {$x$};
  \draw[-Latex, grigio] (0,-0.3) -- (0,3.6) node[above, black] {$y$};
  \draw[blu, dashed] (1.5,0.8) rectangle (5,3);
  \foreach \p/\lab/\pos in {(1.5,0.8)/{(x_{0},y_{0})}/below left,
      (5,0.8)/{(x,y_{0})}/below right, (1.5,3)/{(x_{0},y)}/above left,
      (5,3)/{(x,y)}/above right}
    {\fill[blu] \p circle (1.5pt); \node[font=\scriptsize, \pos] at \p {$\lab$};}
  \fill[rame] (3.9,1.6) circle (1.5pt) node[right, font=\scriptsize] {$(x_{1},y_{1})$};
  \fill[verde] (2.4,2.3) circle (1.5pt) node[right, font=\scriptsize] {$(x_{2},y_{2})$};
  \draw[tratteggio] (1.5,0) node[below, black, font=\scriptsize] {$x_{0}$} -- (1.5,0.8);
  \draw[tratteggio] (5,0) node[below, black, font=\scriptsize] {$x$} -- (5,0.8);
  \draw[tratteggio] (0,0.8) node[left, black, font=\scriptsize] {$y_{0}$} -- (1.5,0.8);
  \draw[tratteggio] (0,3) node[left, black, font=\scriptsize] {$y$} -- (1.5,3);
\end{tikzpicture}
\caption{I punti usati nella dimostrazione del teorema di Schwarz: i quattro vertici del
rettangolo, nei quali si calcola $f$, e i due punti interni $(x_{1},y_{1})$ e
$(x_{2},y_{2})$ forniti dal teorema di Lagrange, nei quali le derivate miste coincidono.}
\label{fig:schwarz}
\end{figure}

\begin{osservazione}[Le ipotesi di continuità non si possono omettere]
\label{oss:schwarz-controesempio}
Se le derivate seconde miste non sono continue, non necessariamente sono uguali.
Consideriamo la funzione
\begin{equation}\label{eq:schwarz-controesempio}
  g(x,y)=\frac{x^{3}y}{x^{2}+y^{2}} \quad\text{per }(x,y)\neq(0,0), \qquad g(0,0)=0 .
\end{equation}
La funzione è nulla sugli assi, quindi $g_{x}(0,0)=g_{y}(0,0)=0$. Per $(x,y)\neq(0,0)$,
\[
  g_{x}=\frac{3x^{2}y(x^{2}+y^{2})-2x^{4}y}{(x^{2}+y^{2})^{2}}, \qquad
  g_{y}=\frac{x^{3}(x^{2}+y^{2})-2x^{3}y^{2}}{(x^{2}+y^{2})^{2}} .
\]
In particolare $g_{x}(0,y)=0$ per ogni $y$, e $g_{y}(x,0)=x$ per ogni $x$. Quindi
\begin{align*}
  g_{xy}(0,0)&=\lim_{k\to0}\frac{g_{x}(0,k)-g_{x}(0,0)}{k}=0,\\
  g_{yx}(0,0)&=\lim_{h\to0}\frac{g_{y}(h,0)-g_{y}(0,0)}{h}=1 :
\end{align*}
le due derivate miste nell'origine sono diverse. Il teorema di Schwarz non è
contraddetto: le derivate miste di $g$ esistono in ogni punto, ma non sono continue
nell'origine.
\end{osservazione}

\begin{python}[Il teorema di Schwarz e il suo controesempio]
import sympy as sp

x, y, h = sp.symbols("x y h", real=True)
f = x**y                                 # definita per x > 0
fxy = sp.diff(f, x, y)              # prima rispetto a x, poi a y
fyx = sp.diff(f, y, x)              # prima rispetto a y, poi a x
print("f_xy =", sp.simplify(fxy), "   f_yx =", sp.simplify(fyx))
print("f_xy - f_yx =", sp.simplify(fxy - fyx))

# il controesempio: g = x^3 y / (x^2 + y^2), g(0, 0) = 0
g = x**3 * y / (x**2 + y**2)
gx0y = sp.simplify(sp.diff(g, x).subs(x, 0))   # g_x(0, y), y != 0
gyx0 = sp.simplify(sp.diff(g, y).subs(y, 0))   # g_y(x, 0), x != 0
print("g_x(0, y) =", gx0y, "   g_y(x, 0) =", gyx0)
# g_x(0,0) = g_y(0,0) = 0 (g è nulla sugli assi), quindi:
print("g_xy(0,0) =", sp.limit(gx0y.subs(y, h) / h, h, 0),
      "   g_yx(0,0) =", sp.limit(gyx0.subs(x, h) / h, h, 0))
\end{python}
\begin{risultato}
f_xy = x**(y - 1)*(y*log(x) + 1)    f_yx = x**(y - 1)*(y*log(x) + 1)
f_xy - f_yx = 0
g_x(0, y) = 0    g_y(x, 0) = x
g_xy(0,0) = 0    g_yx(0,0) = 1
\end{risultato}

Il teorema di Schwarz fornisce informazioni anche sulle derivate di ordine superiore:
applicandolo, ad esempio, alla derivata parziale $f_{x}$, si trova $f_{xxy}=f_{xyx}$, purché
tali derivate terze siano continue. Pertanto, se una funzione di due variabili ha
derivate parziali continue fino al terzo ordine, le derivate terze distinte sono al più
quattro: $f_{xxx}$, $f_{xxy}$,
$f_{xyy}$, $f_{yyy}$. Conta soltanto quante volte si deriva rispetto a ciascuna variabile,
non l'ordine.

\begin{idea}[L'equazione di Laplace e le correnti potenziali]
In aerodinamica una corrente piana incompressibile e irrotazionale si descrive con un
\emph{potenziale di velocità} $\varphi(x,y)$, le cui derivate parziali sono le componenti
della velocità: $u=\varphi_{x}$, $v=\varphi_{y}$. La conservazione della massa impone
$u_{x}+v_{y}=0$, cioè
\begin{equation}\label{eq:laplace}
  \varphi_{xx}+\varphi_{yy}=0 ,
\end{equation}
l'\emph{equazione di Laplace}; le sue soluzioni si dicono \emph{funzioni armoniche}.
L'irrotazionalità, $v_{x}=u_{y}$, è invece automatica: è proprio il teorema di Schwarz,
$\varphi_{yx}=\varphi_{xy}$. Sono armoniche, per $(x,y)\neq(0,0)$, il potenziale di una
\emph{sorgente} $\log\sqrt{x^{2}+y^{2}}$ (che è, a meno del fattore $\frac12$, la prima delle funzioni
\eqref{eq:radiali}), quello di un \emph{vortice} $\operatorname{arctg}\frac yx$ (nei punti
con $x\neq0$) e quello di una \emph{doppietta} $\frac{x}{x^{2}+y^{2}}$. Sommando la corrente uniforme $\varphi=x$ e la
doppietta si ottiene la corrente attorno a un cilindro circolare: è il primo passo della
teoria classica dei profili alari, che conduce al teorema di Kutta--Žukovskij sulla
portanza. Il programma verifica l'equazione di Laplace con \texttt{sympy}.
\end{idea}

\begin{python}[Funzioni armoniche delle correnti potenziali]
import sympy as sp

x, y = sp.symbols("x y", real=True)
lap = lambda u: sp.simplify(sp.diff(u, x, 2) + sp.diff(u, y, 2))
for nome, phi in (("sorgente", sp.log(sp.sqrt(x**2 + y**2))),
                  ("vortice", sp.atan(y / x)),
                  ("doppietta", x / (x**2 + y**2)),
                  ("cilindro", x + x / (x**2 + y**2)),
                  ("non armonica", x**2 + y**2)):
    print(f"{nome:13} phi = {str(phi):24}  laplaciano = {lap(phi)}")
\end{python}
\begin{risultato}
sorgente      phi = log(sqrt(x**2 + y**2))    laplaciano = 0
vortice       phi = atan(y/x)                 laplaciano = 0
doppietta     phi = x/(x**2 + y**2)           laplaciano = 0
cilindro      phi = x + x/(x**2 + y**2)       laplaciano = 0
non armonica  phi = x**2 + y**2               laplaciano = 4
\end{risultato}

\begin{notastorica}[Eulero, Clairaut, Schwarz, Peano]
L'uguaglianza delle derivate miste fu data per scontata da Nicolaus Bernoulli già nel
1721, e usata da Leonhard Euler in una memoria sulle famiglie di curve scritta nel 1734 (e
pubblicata nel 1740). Alexis Clairaut, nelle ricerche sul calcolo integrale del
1739--1740, se ne servì per un problema che abbiamo appena incontrato: riconoscere quando
un'espressione $M\,dx+N\,dy$ è il differenziale di una funzione $\varphi$. Se lo è, si ha
$M=\varphi_{x}$ e $N=\varphi_{y}$, e l'uguaglianza delle derivate miste impone
$M_{y}=N_{x}$; per una corrente piana, con $M=u$ e $N=v$, è la condizione di
irrotazionalità $u_{y}=v_{x}$. Eulero raccolse questi risultati, insieme al teorema sulle
funzioni omogenee (\cref{subsec:funzioni-omogenee}), nel capitolo sulle funzioni di più
variabili delle \emph{Institutiones calculi differentialis} (1755). Le loro dimostrazioni,
e quelle che seguirono per oltre un secolo, da Lagrange (1797) a Cauchy (1823), erano
incomplete. Nel 1867 Lorenz Leonard Lindelöf,
analizzando sistematicamente le dimostrazioni precedenti, mostrò con un controesempio che
una qualche ipotesi è necessaria; la prima dimostrazione rigorosa, con ipotesi di
continuità del tipo di quelle del \cref{teo:schwarz}, fu data da Hermann Amandus Schwarz
nel 1873. Il controesempio oggi classico, $\frac{xy(x^{2}-y^{2})}{x^{2}+y^{2}}$, dello stesso
tipo di quello dell'\cref{oss:schwarz-controesempio}, fu presentato da Giuseppe Peano nel
1884, nelle note al \emph{Calcolo differenziale} di Genocchi.
\end{notastorica}

\subsection{Approfondimento: un primo sguardo alle equazioni alle derivate parziali}
\label{subsec:edp-primo-sguardo}
\index{equazione!alle derivate parziali}

L'equazione di Laplace \eqref{eq:laplace} è il primo esempio che incontriamo di
\emph{equazione alle derivate parziali}: una relazione tra una funzione incognita di più
variabili e le sue derivate parziali, che si dice di ordine $k$ se vi compaiono derivate
fino all'ordine $k$. Accanto a quella di Laplace, la fisica matematica ha altre due
equazioni modello del secondo ordine, nelle quali una delle variabili è il tempo $t$:
\begin{align}
  u_{t}&=k\,u_{xx} &&\text{(equazione del calore)},\label{eq:calore}\\
  u_{tt}&=c^{2}u_{xx} &&\text{(equazione delle onde)},\label{eq:onde}
\end{align}
con $k>0$ e $c>0$ costanti. La prima descrive la temperatura $u(x,t)$ nel punto di
ascissa $x$ di una sbarra omogenea, all'istante $t$; la costante $k$ è la diffusività
termica del materiale. La seconda descrive lo spostamento $u(x,t)$ di una corda vibrante,
oppure la perturbazione di pressione di un'onda sonora in un condotto; vedremo che $c$ è
la velocità con cui le perturbazioni si propagano. Con gli strumenti di questo capitolo
possiamo verificare che alcune funzioni esplicite le risolvono, e capire da esse
\emph{che cosa} le equazioni descrivono.

\needspace{9\baselineskip}
\begin{proposizione}[Soluzione di d'Alembert dell'equazione delle onde]\label{prop:dalembert}
\index{equazione!delle onde}
Siano $F$ e $G$ funzioni di una variabile derivabili due volte in $\R$, e sia $c>0$.
Allora la funzione
\begin{equation}\label{eq:dalembert}
  u(x,t)=F(x-ct)+G(x+ct)
\end{equation}
soddisfa l'equazione delle onde \eqref{eq:onde} in ogni punto $(x,t)\in\R^{2}$.
\end{proposizione}

\begin{schema}
\Ipotesi{$F$, $G$ derivabili due volte in $\R$; \ $c>0$.}
\Tesi{$u(x,t)=F(x-ct)+G(x+ct)$ verifica $u_{tt}=c^{2}u_{xx}$ in $\R^{2}$.}
\end{schema}

\begin{strategia}
Ogni derivata parziale è la derivata di una funzione di una variabile composta con la
funzione lineare $x-ct$ (oppure $x+ct$): basta la regola della catena del
\cref{teo:derivata-composta}. La funzione interna ha derivata $1$ rispetto a $x$ e $\mp c$
rispetto a $t$; derivando due volte, il fattore $\mp c$ compare al quadrato e produce il
$c^{2}$ dell'equazione.
\end{strategia}

\begin{dimostrazione}
Fissato $t$, la funzione $x\mapsto F(x-ct)$ è la composta di $F$ con la funzione
$x\mapsto x-ct$, che ha derivata $1$; per il \cref{teo:derivata-composta}
$\frac{\partial}{\partial x}F(x-ct)=F'(x-ct)$ e, derivando ancora,
$\frac{\partial^{2}}{\partial x^{2}}F(x-ct)=F''(x-ct)$. Fissato $x$, la funzione
$t\mapsto x-ct$ ha derivata $-c$, quindi
\[
  \frac{\partial}{\partial t}F(x-ct)=-c\,F'(x-ct), \qquad
  \frac{\partial^{2}}{\partial t^{2}}F(x-ct)=c^{2}F''(x-ct).
\]
Allo stesso modo, con $+c$ al posto di $-c$, le derivate seconde di $G(x+ct)$ rispetto a
$x$ e a $t$ valgono $G''(x+ct)$ e $c^{2}G''(x+ct)$. Sommando,
$u_{tt}=c^{2}\bigl[F''(x-ct)+G''(x+ct)\bigr]=c^{2}u_{xx}$.
\end{dimostrazione}

Il significato della \eqref{eq:dalembert} è trasparente: il grafico di $x\mapsto F(x-ct)$
è quello di $F$ traslato di $ct$ verso destra, quindi $F(x-ct)$ è un'onda che viaggia
verso destra con velocità $c$ senza deformarsi, e $G(x+ct)$ un'onda che viaggia verso
sinistra. Ad esempio, $u=\frac12\bigl[F(x-ct)+F(x+ct)\bigr]$ vale $F(x)$ per $t=0$ e ha
$u_{t}(x,0)=\frac12\bigl[-cF'(x)+cF'(x)\bigr]=0$: è la corda che, deformata con il profilo $F$,
viene lasciata andare da ferma. Il profilo si divide in due metà uguali, che si
allontanano in versi opposti (\cref{fig:edp-onde-calore}, a sinistra).

\begin{figure}[htb]
\centering
\includegraphics[width=\textwidth]{edp_onde_calore}
\caption{A sinistra, la soluzione $u=\frac12[F(x-ct)+F(x+ct)]$ dell'equazione delle onde,
con $F(x)=e^{-4x^{2}}$ e $c=1$: il profilo iniziale si divide in due onde che viaggiano
senza deformarsi in versi opposti. A destra, la soluzione \eqref{eq:gaussiana-calore}
dell'equazione del calore, con $k=1$: la campana si abbassa e si allarga
(\texttt{python/fig\_edp\_onde\_calore.py}).}
\label{fig:edp-onde-calore}
\end{figure}

\index{equazione!del calore}Per l'equazione del calore la verifica è altrettanto diretta. La funzione
$u=e^{-kn^{2}t}\sen(nx)$, con $n\in\N$, ha $u_{t}=-kn^{2}u$ e $u_{xx}=-n^{2}u$, e quindi
soddisfa la \eqref{eq:calore}: un profilo sinusoidale di temperatura si smorza
esponenzialmente, tanto più in fretta quanto più è fitto. Più istruttiva è la funzione,
definita per $t>0$,
\begin{equation}\label{eq:gaussiana-calore}
  u(x,t)=\frac{1}{\sqrt{4\pi kt}}\,e^{-x^{2}/(4kt)},
\end{equation}
che il programma più avanti verifica. Per $t$ piccolo è una campana alta e stretta: il
calore depositato in un punto della sbarra all'istante iniziale. Col tempo la campana si
abbassa e si allarga, con una larghezza proporzionale a $\sqrt{kt}$
(\cref{fig:edp-onde-calore}, a destra): là un profilo viaggia senza deformarsi, qui si
smussa e si diffonde.

\begin{esempio}[Le onde di Mach]\label{es:onde-mach}
Un profilo alare sottile sia investito da una corrente supersonica, con numero di Mach
$M>1$. Se le perturbazioni prodotte dal profilo sono piccole, il loro potenziale di
velocità $\varphi(x,y)$, con l'asse $x$ nella direzione della corrente, soddisfa l'equazione
linearizzata
\begin{equation}\label{eq:mach-linearizzata}
  (M^{2}-1)\,\varphi_{xx}-\varphi_{yy}=0 .
\end{equation}
Posto $\beta=\sqrt{M^{2}-1}$, è l'equazione delle onde \eqref{eq:onde} con $y$ al posto di $t$
e $\beta$ al posto di $c$; per la \cref{prop:dalembert} è risolta da
$\varphi=F(x-\beta y)+G(x+\beta y)$. Sopra il profilo, dove le perturbazioni sono portate a
valle, si ha $\varphi=F(x-\beta y)$: la perturbazione è costante lungo le rette
$x-\beta y=\text{costante}$, inclinate sulla direzione della corrente di un angolo $\mu$ con
$\tg\mu=\frac1\beta$, cioè $\sen\mu=\frac1M$. Sono le \emph{onde di Mach}, e $\mu$ è
l'\emph{angolo di Mach}: per $M=2$ si ha $\mu=30^{\circ}$. In regime subsonico, $M<1$,
l'equazione diventa $(1-M^{2})\varphi_{xx}+\varphi_{yy}=0$, e la sostituzione
$\varphi(x,y)=\Phi\bigl(x,\sqrt{1-M^{2}}\,y\bigr)$ la riconduce all'equazione di Laplace
$\Phi_{xx}+\Phi_{yy}=0$ (è la base della regola di Prandtl--Glauert): non ci sono più
direzioni privilegiate, e il profilo «si fa sentire» in tutto il campo, anche a monte.
\end{esempio}

\needspace{7\baselineskip}
\begin{python}[Soluzioni esplicite delle equazioni delle onde e del calore]
import math
import sympy as sp

x, y, t, n = sp.symbols("x y t n", real=True)
k, c, M = sp.symbols("k c M", positive=True)
F, G, d = sp.Function("F"), sp.Function("G"), sp.diff
onde = lambda u: sp.simplify(d(u, t, 2) - c**2 * d(u, x, 2))
calore = lambda u: sp.simplify(d(u, t) - k * d(u, x, 2))
print("u_tt - c^2 u_xx per F(x-ct)+G(x+ct):",
      onde(F(x - c*t) + G(x + c*t)))
print("u_t - k u_xx per e^(-kn^2 t)sen(nx):",
      calore(sp.exp(-k * n**2 * t) * sp.sin(n*x)))
gauss = sp.exp(-x**2 / (4*k*t)) / sp.sqrt(4 * sp.pi * k * t)
print("u_t - k u_xx per la gaussiana:      ", calore(gauss))
# onde di Mach: (M^2 - 1) phi_xx - phi_yy = 0 per phi = F(x - beta y)
beta = sp.sqrt(M**2 - 1)
phi = F(x - beta * y)
print("(M^2-1) phi_xx - phi_yy:",
      sp.simplify((M**2 - 1) * d(phi, x, 2) - d(phi, y, 2)))
for Mach in (1.2, 1.5, 2.0, 3.0):
    mu = math.degrees(math.asin(1 / Mach))
    print(f"M = {Mach:.1f}:  angolo di Mach = {mu:5.2f} gradi")
\end{python}
\begin{risultato}
u_tt - c^2 u_xx per F(x-ct)+G(x+ct): 0
u_t - k u_xx per e^(-kn^2 t)sen(nx): 0
u_t - k u_xx per la gaussiana:       0
(M^2-1) phi_xx - phi_yy: 0
M = 1.2:  angolo di Mach = 56.44 gradi
M = 1.5:  angolo di Mach = 41.81 gradi
M = 2.0:  angolo di Mach = 30.00 gradi
M = 3.0:  angolo di Mach = 19.47 gradi
\end{risultato}

Al crescere di $M$ l'angolo di Mach diminuisce e le onde si ripiegano all'indietro.

\begin{osservazione}[Tre equazioni modello e la fluidodinamica computazionale]
\label{oss:edp-cfd}
Le equazioni di Laplace, del calore e delle onde sono i prototipi di tre famiglie, dette
rispettivamente \emph{ellittiche}, \emph{paraboliche} e \emph{iperboliche}, dai
comportamenti molto diversi. Nelle prime ogni punto risente di tutto il contorno: è il
caso delle correnti subsoniche a potenziale, che i \emph{metodi a pannelli} calcolano
distribuendo sulla superficie del velivolo sorgenti, vortici e doppiette come quelli
dell'idea «L'equazione di Laplace e le correnti potenziali». Le seconde descrivono la
diffusione: nei fluidi è la quantità di moto a diffondere, con la viscosità cinematica
$\nu$ al posto di $k$, e lo spessore della zona rallentata vicino a una parete cresce come
$\sqrt{\nu t}$, come la larghezza della campana \eqref{eq:gaussiana-calore}; è il primo
passo verso la teoria dello strato limite. Le terze descrivono propagazioni a velocità
finita lungo direzioni privilegiate, come l'acustica e le correnti supersoniche
dell'\cref{es:onde-mach}. Le equazioni di Navier--Stokes, risolte dai codici di CFD,
mescolano i tre comportamenti, e un buon metodo numerico rispetta la natura di ciascuno:
nelle regioni iperboliche, ad esempio, gli schemi \emph{upwind} prendono l'informazione
dalla direzione da cui arrivano le onde.
\end{osservazione}

\begin{notastorica}[Le equazioni della fisica matematica]
Le equazioni alle derivate parziali compaiono, ancora sporadiche, nella prima metà del
Settecento. Nella stessa memoria del 1747 in cui scrisse l'equazione della corda vibrante
(si veda la nota storica del \S\ref{sec:derivate-parziali}), d'Alembert ne diede la
soluzione come somma di due onde che viaggiano in versi opposti, cioè la
\cref{prop:dalembert}; il dibattito che ne seguì con Eulero e con Daniel Bernoulli su quali
funzioni $F$ e $G$ fossero ammissibili pesò a lungo sul concetto stesso di funzione
(nota storica del \S\ref{sec:continue}). L'equazione di
Laplace nacque invece dai fluidi: Eulero la scrisse per il potenziale di velocità di una
corrente incomprimibile nei \emph{Principia motus fluidorum}, presentati all'Accademia di
Berlino nel 1752; Pierre-Simon Laplace la ritrovò nel 1782 per il potenziale
gravitazionale, e Siméon-Denis Poisson, nel 1813, mostrò che dentro la materia va corretta
con un secondo membro proporzionale alla densità. Con Fourier, che nel 1807 presentò
l'equazione del calore, con Laplace e con Poisson, la scuola francese fece delle equazioni
alle derivate parziali uno dei campi principali dell'analisi; a Cauchy e a Sof'ja
Kovalevskaja (1874) si deve un teorema di esistenza delle soluzioni. Riemann,
infine, osservò che una funzione armonica con valori assegnati sul bordo di un dominio si
ottiene minimizzando un integrale (il «principio di Dirichlet»); ma che il minimo esista
non è affatto ovvio, e lo dimostrò soltanto David Hilbert, nel 1900.
\end{notastorica}

%% ----------------------------------------------------------------------------
%%  §7.5  Massimi e minimi relativi
%% ----------------------------------------------------------------------------
\section{Massimi e minimi relativi}
\label{sec:estremi-due}

Le definizioni di massimo e di minimo relativo si estendono senza difficoltà,
sostituendo gli intervalli con gli intorni circolari.

\begin{definizione}[Massimi e minimi relativi]\label{def:estremi-due}
\index{massimo!relativo!in più variabili}\index{minimo!relativo!in più variabili}
Sia $f$ una funzione di due variabili definita in un insieme $D\subseteq\R^{2}$. Un punto
$(x_{0},y_{0})\in D$ è di \emph{massimo relativo} per $f$ in $D$ se esiste un intorno
circolare $I_{\delta}$ di $(x_{0},y_{0})$ tale che
\begin{equation}\label{eq:max-due}
  f(x_{0},y_{0})\ge f(x,y) \qquad\perogni(x,y)\in I_{\delta}\cap D ;
\end{equation}
è di \emph{minimo relativo} se $f(x_{0},y_{0})\le f(x,y)$ per ogni
$(x,y)\in I_{\delta}\cap D$. Se le disuguaglianze valgono in senso stretto per
$(x,y)\neq(x_{0},y_{0})$, il massimo o il minimo si dice \emph{stretto}.
\end{definizione}

Come per le funzioni di una variabile, la derivata dà informazioni soltanto nei punti
«interni».

\begin{definizione}[Punto interno]\label{def:punto-interno}\index{punto!interno}
Un punto $(x_{0},y_{0})$ è \emph{interno} all'insieme $D\subseteq\R^{2}$ se esiste un
intorno circolare $I_{\delta}$ di $(x_{0},y_{0})$ contenuto in $D$. Un insieme i cui punti
sono tutti interni si dice \emph{aperto}.
\end{definizione}

La \cref{fig:punto-interno} illustra la differenza tra i punti interni e quelli che
stanno sul bordo dell'insieme.

\begin{figure}[htb]
\centering
\begin{tikzpicture}[scale=1.0]
  \draw[-Latex, grigio] (-0.3,0) -- (6.4,0) node[right, black] {$x$};
  \draw[-Latex, grigio] (0,-0.3) -- (0,3.6) node[above, black] {$y$};
  \fill[pattern=north east lines, pattern color=blu]
    (0.8,3.2) .. controls (0.6,1.2) and (2.0,0.1) .. (3.5,0.6)
    .. controls (5.0,1.1) and (6.0,2.4) .. (5.9,3.3) -- (0.8,3.2) -- cycle;
  \draw[blu, thick]
    (0.8,3.2) .. controls (0.6,1.2) and (2.0,0.1) .. (3.5,0.6)
    .. controls (5.0,1.1) and (6.0,2.4) .. (5.9,3.3);
  \fill[white] (2.2,2.0) circle (0.75);
  \fill[pattern=north east lines, pattern color=blu!50] (2.2,2.0) circle (0.75);
  \draw[rame, thick, dashed] (2.2,2.0) circle (0.75);
  \fill (2.2,2.0) circle (1.5pt) node[above right, font=\scriptsize] {$(x_{0},y_{0})$};
  \node[font=\scriptsize, fill=white, inner sep=1pt] at (2.2,1.6) {punto interno};
  \fill (5.0216,1.507) circle (1.5pt); % punto della curva di Bézier (t = 0,4)
  \node[font=\scriptsize, right] at (5.25,1.25) {punto non interno};
  \node[font=\scriptsize, blu] at (4.6,2.6) {$D$};
\end{tikzpicture}
\caption{Un punto interno all'insieme $D$, che ha un intorno circolare tutto contenuto in
$D$, e un punto di $D$ che sta sul bordo, ogni intorno del quale contiene punti esterni a
$D$.}
\label{fig:punto-interno}
\end{figure}

\begin{teorema}[Condizione necessaria per gli estremi]\label{teo:fermat-due}
\index{teorema!di Fermat!in più variabili}
Se $(x_{0},y_{0})$ è un punto di massimo o di minimo relativo interno al dominio $D$ di una
funzione $f(x,y)$, e se $f$ è dotata di derivate parziali prime in $(x_{0},y_{0})$, allora
\begin{equation}\label{eq:fermat-due}
  f_{x}(x_{0},y_{0})=0, \qquad f_{y}(x_{0},y_{0})=0 .
\end{equation}
\end{teorema}

\begin{schema}
\Ipotesi{$(x_{0},y_{0})$ interno a $D$, di massimo (o minimo) relativo; esistono
$f_{x}(x_{0},y_{0})$, $f_{y}(x_{0},y_{0})$.}
\Tesi{$\nabla f(x_{0},y_{0})=(0,0)$.}
\end{schema}

\begin{strategia}
Restringiamo $f$ alle rette per $(x_{0},y_{0})$ parallele agli assi. Si ottengono funzioni
di una variabile che hanno un estremo relativo interno, e alle quali si applica il
teorema di Fermat; le loro derivate sono proprio le derivate parziali.
\end{strategia}

\begin{dimostrazione}
Supponiamo che $(x_{0},y_{0})$ sia un punto di massimo relativo interno; esiste allora un
intorno circolare $I_{\delta}$, contenuto in $D$, nel quale vale la \eqref{eq:max-due}.
Fissato $y=y_{0}$, consideriamo la funzione di una variabile reale
$F(x)=f(x,y_{0})$, definita nell'intervallo $(x_{0}-\delta,x_{0}+\delta)$ (i punti
$(x,y_{0})$ con $\abs{x-x_{0}}<\delta$ stanno in $I_{\delta}$). Per la \eqref{eq:max-due},
$F(x_{0})\ge F(x)$ per ogni $x$ di tale intervallo: $x_{0}$ è un punto di massimo
relativo interno per $F$. Per il teorema di Fermat (\cref{teo:fermat}),
$F'(x_{0})=f_{x}(x_{0},y_{0})=0$. Si procede in modo analogo, con la funzione
$G(y)=f(x_{0},y)$, per ottenere $f_{y}(x_{0},y_{0})=0$.
\end{dimostrazione}

I punti di massimo o di minimo interni al dominio vanno dunque ricercati tra i punti che
annullano entrambe le derivate parziali prime; tali punti si dicono \emph{punti critici}
(o \emph{stazionari}) della funzione: sono i punti in cui il gradiente si annulla. Come
per le funzioni di una variabile, la condizione è necessaria ma non sufficiente: la
funzione $y^{2}-x^{2}$ ha gradiente $(-2x,2y)$, nullo nell'origine, ma l'origine non è né
di massimo né di minimo, perché $f(x,0)=-x^{2}<0$ e $f(0,y)=y^{2}>0$ per $x,y\neq0$. Un
punto critico di questo tipo si chiama \emph{punto di sella}, per la forma del grafico
(\cref{fig:superfici}).

\subsection{Il criterio dell'hessiano}

Il teorema che segue fornisce un criterio per stabilire la natura di un punto critico.
Supponiamo che $f$ abbia derivate seconde continue in un insieme aperto $D$; si definisce
il \emph{determinante hessiano} (o semplicemente \emph{hessiano})
\begin{equation}\label{eq:hessiano}
  H(x,y)=\begin{vmatrix} f_{xx}(x,y) & f_{xy}(x,y)\\ f_{yx}(x,y) & f_{yy}(x,y)\end{vmatrix}
  =f_{xx}(x,y)\,f_{yy}(x,y)-\bigl(f_{xy}(x,y)\bigr)^{2},
\end{equation}
dove si è tenuto conto che, per il teorema di Schwarz, $f_{xy}=f_{yx}$. Talvolta, per
mettere in luce la dipendenza dalla funzione $f$, si scrive $Hf(x,y)$.

\begin{teorema}[Criterio dell'hessiano]\label{teo:hessiano}\index{criterio!dell'hessiano}
\index{hessiano}
Sia $f$ dotata di derivate parziali prime e seconde continue in un intorno del punto
$(x_{0},y_{0})$, e sia $(x_{0},y_{0})$ un punto critico:
$f_{x}(x_{0},y_{0})=f_{y}(x_{0},y_{0})=0$.
\begin{enumerate}[label=\textup{(\roman*)}]
  \item Se $Hf(x_{0},y_{0})>0$ e $f_{xx}(x_{0},y_{0})>0$, allora $(x_{0},y_{0})$ è un punto di
  minimo relativo (stretto).
  \item Se $Hf(x_{0},y_{0})>0$ e $f_{xx}(x_{0},y_{0})<0$, allora $(x_{0},y_{0})$ è un punto di
  massimo relativo (stretto).
  \item Se $Hf(x_{0},y_{0})<0$, allora $(x_{0},y_{0})$ non è né di massimo né di minimo
  relativo: si dice che è un \emph{punto di sella}.
\end{enumerate}
\end{teorema}

\begin{schema}
\Ipotesi{$f$ con derivate prime e seconde continue vicino a $(x_{0},y_{0})$;
$\nabla f(x_{0},y_{0})=(0,0)$.}
\Tesi{$H>0$, $f_{xx}>0$: minimo; \ $H>0$, $f_{xx}<0$: massimo; \ $H<0$: sella (tutto
calcolato in $(x_{0},y_{0})$).}
\end{schema}

\begin{strategia}
Restringiamo $f$ al segmento che congiunge $(x_{0},y_{0})$ con un punto vicino
$(x_{0}+h,y_{0}+k)$: si ottiene una funzione di una variabile $g(t)$ con $g'(0)=0$, la cui
derivata seconda è la \emph{forma quadratica}
$f_{xx}h^{2}+2f_{xy}hk+f_{yy}k^{2}$. Il segno di questa forma, per tutte le direzioni $(h,k)$,
è deciso proprio da $f_{xx}$ e da $H$ (completando il quadrato). Per rendere rigoroso il
ragionamento occorrono due strumenti: la derivata di $f$ lungo un segmento, che richiede
la nozione di differenziabilità, e una formula di Taylor con il resto nella forma di
Lagrange. Per questo la dimostrazione è rinviata al \cref{sec:dim-hessiano}, in
appendice.
\end{strategia}

Un punto critico è dunque di massimo o di minimo relativo se $Hf(x_{0},y_{0})>0$, mentre
se $Hf(x_{0},y_{0})<0$ non è né di massimo né di minimo. Il caso $Hf(x_{0},y_{0})=0$ non è
contemplato nell'enunciato, perché in tal caso il punto può essere di massimo, di minimo
o di sella: nell'origine le funzioni $x^{4}+y^{4}$, $-(x^{4}+y^{4})$ e
$x^{4}-y^{4}$ hanno tutte $H=0$, ma la prima ha un minimo, la seconda un massimo, la terza
una sella. Si noti infine che, se $H>0$, allora $f_{xx}f_{yy}>f_{xy}^{2}\ge0$, e quindi
$f_{xx}$ e $f_{yy}$ hanno lo stesso segno: nel criterio si può usare l'una o l'altra.

\begin{esempio}[Quattro funzioni e i loro punti critici]\label{es:punti-critici}
\begin{enumerate}[label=\textup{(\alph*)}]
  \item La funzione $f(x,y)=x^{2}+y^{2}$ ha derivate parziali $f_{x}=2x$, $f_{y}=2y$, che si
  annullano nel punto $(0,0)$. L'hessiano vale
  $H=\begin{vmatrix}2&0\\0&2\end{vmatrix}=4>0$ ed $f_{xx}=2>0$: l'origine è di minimo
  relativo. Si verifica anche direttamente che è di minimo \emph{assoluto}, perché
  $f(0,0)=0\le x^{2}+y^{2}$ per ogni $(x,y)$.
  \item La funzione $f(x,y)=y^{2}-x^{2}$ ha l'unico punto critico $(0,0)$, con
  $H=\begin{vmatrix}-2&0\\0&2\end{vmatrix}=-4<0$: è un punto di sella, ben riconoscibile
  nel grafico del paraboloide iperbolico.
  \item La funzione $f(x,y)=x^{3}-6y(x+y)=x^{3}-6xy-6y^{2}$, definita su $\R^{2}$, ha i
  punti critici che risolvono il sistema
  \[
    f_{x}=3x^{2}-6y=0, \qquad f_{y}=-6x-12y=0 .
  \]
  Dalla seconda equazione $x=-2y$, che sostituito nella prima dà $12y^{2}-6y=0$, cioè
  $y=0$ oppure $y=\frac12$; in corrispondenza $x=0$ oppure $x=-1$. I punti critici sono
  $(0,0)$ e $(-1,\frac12)$. Poiché $f_{xx}=6x$, $f_{xy}=-6$, $f_{yy}=-12$,
  \[
    H(x,y)=\begin{vmatrix}6x&-6\\-6&-12\end{vmatrix}=-72x-36 .
  \]
  Nell'origine $H(0,0)=-36<0$: è un punto di sella. Nel punto $(-1,\frac12)$ si ha
  $H=36>0$ ed $f_{xx}(-1,\frac12)=-6<0$: è un punto di massimo relativo, con
  $f(-1,\frac12)=\frac12$ (\cref{fig:punti-critici}, a sinistra). Non è un massimo
  assoluto, perché $f(x,0)=x^{3}\to+\infty$ per $x\to+\infty$.
  \item Per $f(x,y)=y\,e^{-(x^{2}+y^{2})}$ il sistema
  \[
    f_{x}=-2xy\,e^{-(x^{2}+y^{2})}=0, \qquad f_{y}=(1-2y^{2})\,e^{-(x^{2}+y^{2})}=0
  \]
  dà, dalla prima equazione, $x=0$ oppure $y=0$; ma $y=0$ è incompatibile con la seconda,
  che richiede $y=\pm\frac{\sqrt2}2$. I punti critici sono $(0,\pm\frac{\sqrt2}2)$. Il
  programma seguente calcola l'hessiano: in entrambi i punti $H=4e^{-1}>0$, mentre
  $f_{xx}(0,\frac{\sqrt2}2)=-\sqrt2\,e^{-1/2}<0$ e $f_{xx}(0,-\frac{\sqrt2}2)=\sqrt2\,e^{-1/2}>0$:
  il primo è di massimo, il secondo di minimo (\cref{fig:punti-critici}, a destra).
\end{enumerate}
\end{esempio}

\begin{figure}[htb]
\centering
\includegraphics[width=\textwidth]{punti_critici}
\caption{A sinistra, le curve di livello di $x^{3}-6y(x+y)$: intorno al massimo
$(-1,\frac12)$ sono curve chiuse; nella sella $(0,0)$ si incrociano due rami della curva
di livello $0$ (in grigio). A destra, il grafico di $y\,e^{-(x^{2}+y^{2})}$ con il massimo e
il minimo (\texttt{python/fig\_punti\_critici.py}).}
\label{fig:punti-critici}
\end{figure}

\begin{python}[Ricerca e classificazione dei punti critici]
import sympy as sp

x, y = sp.symbols("x y", real=True)

def classifica(f):
    """Punti critici di f(x, y) e loro natura con l'hessiano."""
    fx, fy = sp.diff(f, x), sp.diff(f, y)
    H = sp.hessian(f, (x, y))
    for P in sp.solve([fx, fy], [x, y], dict=True):
        h = sp.simplify(H.det().subs(P))
        fxx = sp.simplify(H[0, 0].subs(P))
        if h > 0:
            tipo = "minimo" if fxx > 0 else "massimo"
        elif h < 0:
            tipo = "sella"
        else:
            tipo = "H = 0: il criterio non decide"
        print(f"({P[x]}, {P[y]}): H = {h}, f_xx = {fxx} -> {tipo}")

for f in (x**2 + y**2, y**2 - x**2, x**3 - 6*y*(x + y),
          y * sp.exp(-(x**2 + y**2)), x**3 + y**3 - 3*x*y):
    print("f =", f)
    classifica(f)
\end{python}
\begin{risultato}
f = x**2 + y**2
(0, 0): H = 4, f_xx = 2 -> minimo
f = -x**2 + y**2
(0, 0): H = -4, f_xx = -2 -> sella
f = x**3 - 6*y*(x + y)
(-1, 1/2): H = 36, f_xx = -6 -> massimo
(0, 0): H = -36, f_xx = 0 -> sella
f = y*exp(-x**2 - y**2)
(0, -sqrt(2)/2): H = 4*exp(-1), f_xx = sqrt(2)*exp(-1/2) -> minimo
(0, sqrt(2)/2): H = 4*exp(-1), f_xx = -sqrt(2)*exp(-1/2) -> massimo
f = x**3 - 3*x*y + y**3
(0, 0): H = -9, f_xx = 0 -> sella
(1, 1): H = 27, f_xx = 6 -> minimo
\end{risultato}

L'ultima funzione dell'elenco, $x^{3}+y^{3}-3xy$, è oggetto di un esercizio del
\cref{sec:esercizi-cap7}. Chiudiamo gli esempi con un punto di massimo in cui
l'hessiano è nullo. La funzione $f(x,y)=\cos(x^{2}+y^{2})$ della \cref{fig:superfici} ha
nell'origine un massimo assoluto, perché $f(0,0)=\cos0=1\ge\cos(x^{2}+y^{2})$ per ogni
$(x,y)$; in accordo con il \cref{teo:fermat-due}, l'origine annulla il gradiente
$\bigl(-2x\sen(x^{2}+y^{2}),\,-2y\sen(x^{2}+y^{2})\bigr)$. Ma tutte le derivate seconde,
\begin{align*}
  f_{xx}&=-2\sen(x^{2}+y^{2})-4x^{2}\cos(x^{2}+y^{2}),\\
  f_{yy}&=-2\sen(x^{2}+y^{2})-4y^{2}\cos(x^{2}+y^{2}),\\
  f_{xy}&=-4xy\cos(x^{2}+y^{2}),
\end{align*}
si annullano nell'origine, e così l'hessiano: il criterio non è applicabile, e la natura
del punto si stabilisce con il confronto diretto.

\subsection{Un'applicazione: il metodo dei minimi quadrati}

In galleria del vento si misura il coefficiente di resistenza $C_{D,i}$ di un modello per
alcuni valori $C_{L,i}$ del coefficiente di portanza, $i=1,\dots,n$. Vogliamo stimare i
parametri della polare parabolica $C_{D}=C_{D0}+kC_{L}^{2}$. Posto $X_{i}=C_{L,i}^{2}$, si
tratta di trovare la retta $C_{D}=b+aX$ che meglio approssima i punti $(X_{i},C_{D,i})$.
Il criterio più usato, dovuto a Legendre e a Gauss, è quello dei \emph{minimi
quadrati}: si scelgono $a$ e $b$ in modo da rendere minima la somma dei quadrati degli
scarti (o \emph{residui})
\begin{equation}\label{eq:minimi-quadrati}
  S(a,b)=\sum_{i=1}^{n}\bigl(C_{D,i}-b-aX_{i}\bigr)^{2} .
\end{equation}
È una funzione di due variabili su tutto $\R^{2}$. I suoi punti critici risolvono il
sistema
\[
  S_{a}=-2\sum_{i=1}^{n}X_{i}\bigl(C_{D,i}-b-aX_{i}\bigr)=0, \qquad
  S_{b}=-2\sum_{i=1}^{n}\bigl(C_{D,i}-b-aX_{i}\bigr)=0,
\]
cioè le \emph{equazioni normali}
\begin{equation}\label{eq:equazioni-normali}
  a\sum_{i}X_{i}^{2}+b\sum_{i}X_{i}=\sum_{i}X_{i}C_{D,i}, \qquad
  a\sum_{i}X_{i}+b\,n=\sum_{i}C_{D,i} .
\end{equation}
È un sistema lineare il cui determinante è $n\sum_{i}X_{i}^{2}-\bigl(\sum_{i}X_{i}\bigr)^{2}$.
Questa quantità è positiva se i valori $X_{i}$ non sono tutti uguali, grazie all'identità
\begin{equation}\label{eq:identita-lagrange}
  n\sum_{i}X_{i}^{2}-\Bigl(\sum_{i}X_{i}\Bigr)^{2}=\frac12\sum_{i=1}^{n}\sum_{j=1}^{n}(X_{i}-X_{j})^{2},
\end{equation}
che si verifica sviluppando il quadrato a secondo membro. Il sistema ha quindi
un'unica soluzione $(a^{*},b^{*})$. Le derivate seconde sono costanti:
$S_{aa}=2\sum_{i}X_{i}^{2}>0$, $S_{ab}=2\sum_{i}X_{i}$, $S_{bb}=2n$, e l'hessiano
$H=4\bigl[n\sum_{i}X_{i}^{2}-(\sum_{i}X_{i})^{2}\bigr]$ è positivo: per il criterio
dell'hessiano $(a^{*},b^{*})$ è un punto di minimo relativo.

In questo caso si può dire di più: il minimo è \emph{assoluto}. Indichiamo con
$r_{i}=C_{D,i}-b^{*}-a^{*}X_{i}$ i residui nel punto di minimo; le equazioni normali dicono
che $\sum_{i}r_{i}=0$ e $\sum_{i}X_{i}r_{i}=0$. Allora, per ogni $(\alpha,\beta)$,
\begin{align*}
  S(a^{*}+\alpha,b^{*}+\beta)&=\sum_{i}\bigl(r_{i}-\alpha X_{i}-\beta\bigr)^{2}\\
  &=S(a^{*},b^{*})-2\alpha\sum_{i}X_{i}r_{i}-2\beta\sum_{i}r_{i}+\sum_{i}(\alpha X_{i}+\beta)^{2},
\end{align*}
e i due termini centrali sono nulli: $S(a^{*}+\alpha,b^{*}+\beta)\ge S(a^{*},b^{*})$, con
l'uguaglianza soltanto per $\alpha=\beta=0$. Il programma applica il metodo a un insieme
di misure simulate.

\begin{python}[Stima di $C_{D0}$ e $k$ con i minimi quadrati]
import numpy as np

# "prove in galleria": C_D misurato per vari C_L (dati simulati)
rng = np.random.default_rng(7)
CL = np.linspace(0.0, 1.2, 13)
CD = 0.025 + 0.0497 * CL**2 + rng.normal(0, 0.0008, CL.size)

# modello C_D = b + a X con X = C_L^2: equazioni normali
X, n = CL**2, CL.size
A = np.array([[np.sum(X**2), np.sum(X)], [np.sum(X), n]])
r = np.array([np.sum(X * CD), np.sum(CD)])
a, b = np.linalg.solve(A, r)
H = 4 * (n * np.sum(X**2) - np.sum(X)**2)   # hessiano di S(a, b)
print(f"k = a = {a:.5f}   C_D0 = b = {b:.5f}   H = {H:.2f} > 0")
print(f"fattore di Oswald (AR = 8): e = {1 / (np.pi * 8 * a):.3f}")
print("controllo con numpy.polyfit:",
      np.round(np.polyfit(X, CD, 1), 5))
S = np.sum((CD - b - a * X)**2)
print(f"somma dei quadrati dei residui S = {S:.3e}")
\end{python}
\begin{risultato}
k = a = 0.05000   C_D0 = b = 0.02478   H = 146.69 > 0
fattore di Oswald (AR = 8): e = 0.796
controllo con numpy.polyfit: [0.05    0.02478]
somma dei quadrati dei residui S = 2.858e-06
\end{risultato}

I dati erano stati simulati con $C_{D0}=0{,}025$ e $k=0{,}0497$ (cioè $e\approx0{,}8$ con
$AR=8$), aggiungendo un rumore casuale di ampiezza tipica $0{,}0008$: il metodo recupera i
parametri con ottima precisione (\cref{fig:minimi-quadrati}). La funzione
\texttt{polyfit} di \texttt{numpy} risolve esattamente lo stesso problema.

\begin{figure}[htb]
\centering
\includegraphics[width=\textwidth]{minimi_quadrati}
\caption{A sinistra, le misure simulate (rame) e la polare parabolica stimata (blu). A
destra, gli stessi dati nel piano $(C_{L}^{2},C_{D})$, dove il modello è una retta. I
residui, cioè i segmenti verticali grigi tra i punti e la retta, la somma dei cui quadrati
è resa minima, sono qui così piccoli (dell'ordine di $10^{-3}$) da restare quasi nascosti
dai punti (\texttt{python/fig\_minimi\_quadrati.py}).}
\label{fig:minimi-quadrati}
\end{figure}

\begin{notastorica}[Gauss, Legendre e l'asteroide perduto]
Il 1º gennaio 1801 Giuseppe Piazzi, dall'osservatorio di Palermo, scoprì Cerere, il primo
asteroide; lo seguì per poche settimane, poi lo perse nel bagliore del Sole. Il
ventiquattrenne Carl Friedrich Gauss, con un metodo di sua invenzione per determinare
l'orbita da poche osservazioni affette da errori, predisse dove ritrovarlo, e alla fine
del 1801 Cerere fu riosservato a meno di mezzo grado dalla posizione prevista. Il metodo
dei minimi quadrati fu pubblicato per la prima volta da Adrien-Marie Legendre nel 1805;
Gauss lo presentò nel 1809, affermando di usarlo dal 1795, e ne diede una giustificazione
probabilistica legata alla curva che porta il suo nome. Le condizioni del secondo ordine
per gli estremi delle funzioni di più variabili risalgono a Lagrange (1759); il
determinante delle derivate seconde porta il nome di Otto Hesse, che lo studiò nel 1844.
\end{notastorica}

%% ----------------------------------------------------------------------------
%%  §7.6  Funzioni di tre o più variabili reali
%% ----------------------------------------------------------------------------
\section{Funzioni di tre o più variabili reali}
\label{sec:n-variabili}

Indichiamo con $\R^{3}$ l'insieme delle \emph{terne ordinate} di numeri reali e, più in
generale, con $\R^{n}$ l'insieme delle \emph{$n$-ple ordinate}:
\begin{equation}\label{eq:Rn}
\begin{gathered}
  \R^{3}=\set{(x,y,z) : x,y,z\in\R},\\
  \R^{n}=\set{(x_{1},x_{2},\dots,x_{n}) : x_{i}\in\R\ \perogni i=1,\dots,n}.
\end{gathered}
\end{equation}
Una legge $f$ che a ogni elemento di un insieme $D\subseteq\R^{n}$ fa corrispondere uno e
un solo numero reale è detta una \emph{funzione di $n$ variabili} e si denota con
$f(x_{1},x_{2},\dots,x_{n})$, oppure, nel caso $n=3$, con $f(x,y,z)$; l'insieme $D$ è il
\emph{dominio} di $f$. La distanza tra due punti di $\R^{n}$ è
$\sqrt{(x_{1}-y_{1})^{2}+\dots+(x_{n}-y_{n})^{2}}$, e con essa si definiscono gli intorni
(sferici, se $n=3$), i limiti e la continuità esattamente come nel piano. Vediamo più in
dettaglio le derivate parziali.

\begin{definizione}[Derivate parziali in $\R^{n}$]\label{def:parziali-n}
La \emph{derivata parziale} di $f$ rispetto alla variabile $x_{i}$ nel punto
$(x_{1},\dots,x_{n})$ è il limite
\begin{equation}\label{eq:parziale-n}
  \lim_{h\to0}\frac{f(x_{1},\dots,x_{i}+h,\dots,x_{n})-f(x_{1},\dots,x_{i},\dots,x_{n})}{h},
\end{equation}
purché esista e sia finito; si denota con uno dei simboli $f_{x_{i}}$,
$\frac{\partial f}{\partial x_{i}}$, $D_{x_{i}}f$, $D_{i}f$. Il vettore
$\nabla f=(f_{x_{1}},\dots,f_{x_{n}})$ è il \emph{gradiente} di $f$.
\end{definizione}

Anche qui si deriva rispetto a una variabile considerando costanti tutte le altre. Ad
esempio, le derivate parziali della funzione di tre variabili
\begin{equation}\label{eq:f-tre-variabili}
  f(x,y,z)=x^{2}+y^{2}+z^{2}+xyz
\end{equation}
sono $f_{x}=2x+yz$, $f_{y}=2y+xz$, $f_{z}=2z+xy$. La funzione di $n$ variabili detta
\emph{modulo} o \emph{norma},
\begin{equation}\label{eq:norma}
  f(x_{1},\dots,x_{n})=\sqrt{x_{1}^{2}+x_{2}^{2}+\dots+x_{n}^{2}},
\end{equation}
che misura la distanza del punto dall'origine, ammette in ogni punto diverso
dall'origine le derivate parziali
\begin{equation}\label{eq:norma-derivate}
  f_{x_{i}}=\frac{2x_{i}}{2\sqrt{x_{1}^{2}+\dots+x_{n}^{2}}}=\frac{x_{i}}{f},
  \qquad i=1,2,\dots,n,
\end{equation}
come si vede derivando la radice con la regola della catena. Il gradiente della norma è
quindi il versore $\frac{(x_{1},\dots,x_{n})}{f}$, diretto radialmente verso l'esterno:
la distanza dall'origine cresce più rapidamente allontanandosi in direzione radiale.

\begin{esempio}[La portanza come funzione di tre variabili]\label{es:portanza-tre}
La portanza $L(\rho,V,C_{L})=\frac12\rho SV^{2}C_{L}$ è una funzione di tre variabili, con
\[
  \frac{\partial L}{\partial\rho}=\frac12SV^{2}C_{L}=\frac L\rho, \qquad
  \frac{\partial L}{\partial V}=\rho SVC_{L}=\frac{2L}V, \qquad
  \frac{\partial L}{\partial C_{L}}=\frac{L}{C_{L}} .
\]
Moltiplicando ciascuna derivata per la variabile corrispondente e dividendo per $L$ si
ottengono le \emph{sensibilità relative} $1$, $2$, $1$: una variazione percentuale piccola
della velocità produce una variazione percentuale \emph{doppia} della portanza, mentre
le variazioni di densità e di $C_{L}$ si trasmettono in egual misura. Ad esempio, salendo
di quota a velocità e assetto costanti, la portanza diminuisce nella stessa proporzione
della densità; per ristabilire l'equilibrio occorre aumentare $C_{L}$ oppure la velocità,
e per quest'ultima basta un aumento percentuale pari a circa la metà. Daremo un
fondamento preciso a questi ragionamenti con il differenziale, nel
\cref{sec:differenziabilita}.
\end{esempio}

Come per le funzioni di due variabili, si definiscono le \emph{derivate parziali
seconde} $f_{x_{i}x_{j}}=(f_{x_{i}})_{x_{j}}$, che vengono a costituire la seguente
matrice quadrata $n\times n$, detta \emph{matrice hessiana}:
\begin{equation}\label{eq:matrice-hessiana}
  \begin{pmatrix}
    f_{x_{1}x_{1}} & f_{x_{1}x_{2}} & \cdots & f_{x_{1}x_{n}}\\
    f_{x_{2}x_{1}} & f_{x_{2}x_{2}} & \cdots & f_{x_{2}x_{n}}\\
    \vdots & \vdots & \ddots & \vdots\\
    f_{x_{n}x_{1}} & f_{x_{n}x_{2}} & \cdots & f_{x_{n}x_{n}}
  \end{pmatrix}.
\end{equation}
Applicando il teorema di Schwarz alla funzione delle sole due variabili $x_{i}$, $x_{j}$
che si ottiene lasciando fissate tutte le altre, con il ruolo di parametri, si ottiene
\begin{equation}\label{eq:schwarz-n}
  f_{x_{i}x_{j}}=f_{x_{j}x_{i}} \qquad\perogni i,j=1,\dots,n,
\end{equation}
purché tali derivate seconde esistano e siano continue. Pertanto la matrice hessiana è
\emph{simmetrica} nel caso in cui tutte le derivate parziali seconde siano continue. Con
le stesse definizioni e la stessa dimostrazione del \cref{teo:fermat-due} si prova che in
un punto di massimo o di minimo relativo interno al dominio risulta
\begin{equation}\label{eq:fermat-n}
  f_{x_{1}}=f_{x_{2}}=\dots=f_{x_{n}}=0,
\end{equation}
nell'ipotesi che le derivate parziali prime esistano nel punto.

\begin{osservazione}[La condizione sufficiente in $n$ variabili]\label{oss:hessiana-n}
Una condizione sufficiente per i massimi e i minimi, in termini della matrice hessiana,
esiste ma è più complessa della formulazione data nel caso di due variabili. Per una
funzione con derivate seconde continue in un intorno del punto critico, essa si
esprime dicendo che la matrice hessiana nel punto critico è \emph{definita positiva}
(minimo) o \emph{definita negativa} (massimo), nozioni che si studiano nei corsi di
Geometria; se la forma quadratica associata assume valori di segno opposto lungo due
direzioni, il punto critico è di sella. Un criterio pratico, dovuto a Sylvester, afferma
che una matrice simmetrica è definita positiva se e solo se sono positivi i minori
principali «di nord-ovest», cioè i determinanti delle sottomatrici formate dalle prime
$k$ righe e colonne, per $k=1,\dots,n$; è definita negativa se e solo se tali minori
hanno segni alterni, a cominciare da un negativo. Per $n=2$ si ritrovano le condizioni
$f_{xx}>0$ (oppure $f_{xx}<0$) e $H>0$ del \cref{teo:hessiano}. Come vedremo nel
\cref{sec:differenziabilita} (per $n=2$; la dimostrazione è identica per $n$ qualsiasi),
la derivata seconda nel punto critico $P$ della restrizione di $f$ alla retta per $P$ di
direzione $v$ è la forma quadratica $v^{T}H(P)\,v$: è questo il motivo per cui, se essa
assume valori di segno opposto per due direzioni, $P$ è un punto di sella.
\end{osservazione}

Il programma seguente trova i punti critici della funzione \eqref{eq:f-tre-variabili} e,
per ciascuno, calcola la forma quadratica $v^{T}Hv$ per i $26$ vettori $v$ con componenti
in $\set{-1,0,1}$ non tutte nulle.

\begin{python}[Punti critici di una funzione di tre variabili]
import itertools
import sympy as sp

x, y, z = sp.symbols("x y z", real=True)
f = x**2 + y**2 + z**2 + x*y*z
grad = [sp.diff(f, v) for v in (x, y, z)]
H = sp.hessian(f, (x, y, z))
print("grad f =", grad)
# direzioni di prova v; la derivata seconda lungo v è v^T H v
dirs = [sp.Matrix(v) for v in itertools.product((-1, 0, 1), repeat=3)
        if any(v)]
for P in sp.solve(grad, [x, y, z], dict=True):
    q = [int((v.T * H.subs(P) * v)[0]) for v in dirs]
    c = ", ".join(f"{int(P[s]):2d}" for s in (x, y, z))
    print(f"P = ({c}):  min v^T H v = {min(q):2d},"
          f"  max v^T H v = {max(q)}")
\end{python}
\begin{risultato}
grad f = [2*x + y*z, x*z + 2*y, x*y + 2*z]
P = (-2, -2, -2):  min v^T H v = -6,  max v^T H v = 10
P = (-2,  2,  2):  min v^T H v = -6,  max v^T H v = 10
P = ( 0,  0,  0):  min v^T H v =  2,  max v^T H v = 6
P = ( 2, -2,  2):  min v^T H v = -6,  max v^T H v = 10
P = ( 2,  2, -2):  min v^T H v = -6,  max v^T H v = 10
\end{risultato}

Nell'origine la matrice hessiana è $2I$, il doppio della matrice identità, e
$v^{T}Hv=2(v_{1}^{2}+v_{2}^{2}+v_{3}^{2})>0$ per ogni $v\neq0$: è definita positiva, e
l'origine è un punto di minimo relativo (come si vede anche osservando che, vicino
all'origine, il termine cubico $xyz$ è trascurabile rispetto a $x^{2}+y^{2}+z^{2}$). Negli
altri quattro punti critici, del tipo $(\pm2,\pm2,\pm2)$ con prodotto delle coordinate
negativo, la forma quadratica assume valori di segno opposto: sono punti di sella.

%% ----------------------------------------------------------------------------
%%  §7.7  Appendice: differenziabilità
%% ----------------------------------------------------------------------------
\section{Appendice: differenziabilità}
\label{sec:differenziabilita}

È ben noto che, per le funzioni di una variabile reale, la derivabilità in un punto
implica la continuità della funzione nel punto (\cref{teo:derivabile-continua}). Non vale
una proprietà analoga per le derivate parziali: esistono funzioni $f(x,y)$, definite in un
intorno di un punto $(x_{0},y_{0})$, che hanno le derivate parziali
$f_{x}(x_{0},y_{0})$, $f_{y}(x_{0},y_{0})$ pur non essendo continue in $(x_{0},y_{0})$.
Consideriamo infatti la funzione
\begin{equation}\label{eq:f-discontinua}
  f(x,y)=\frac{xy}{x^{2}+y^{2}} \quad\text{se }(x,y)\neq(0,0), \qquad f(0,0)=0 .
\end{equation}
La funzione si annulla, oltre che in $(0,0)$, in tutti gli altri punti degli assi
coordinati; pertanto
\[
  f_{x}(0,0)=\lim_{h\to0}\frac{f(h,0)-f(0,0)}{h}=\lim_{h\to0}0=0,
\]
e analogamente $f_{y}(0,0)=0$. Però $f$ non è continua in $(0,0)$, perché non ha limite
per $(x,y)\to(0,0)$: lungo la retta $y=x$ vale costantemente $\frac12$, lungo la retta
$y=-x$ vale $-\frac12$ (\cref{es:xy-su-x2y2,fig:discontinua}). Il motivo è chiaro: le
derivate parziali descrivono il comportamento di $f$ soltanto lungo due rette, e non
vedono che cosa accade nelle altre direzioni. La nozione naturale per le funzioni di più
variabili, che estende davvero il concetto di derivabilità, è quella di
\emph{differenziabilità}.

\begin{definizione}[Funzione differenziabile]\label{def:differenziabile}
\index{funzione!differenziabile}\index{differenziale}
Una funzione $f$, definita in un intorno del punto $(x,y)$, è \emph{differenziabile} in
$(x,y)$ se esistono in tale punto le derivate parziali $f_{x}$, $f_{y}$ e se risulta
\begin{equation}\label{eq:differenziabile}
  \lim_{(h,k)\to(0,0)}\frac{f(x+h,y+k)-f(x,y)-f_{x}(x,y)\,h-f_{y}(x,y)\,k}{\sqrt{h^{2}+k^{2}}}=0 .
\end{equation}
In tal caso la funzione lineare $(h,k)\mapsto f_{x}(x,y)\,h+f_{y}(x,y)\,k$ si chiama il
\emph{differenziale} di $f$ in $(x,y)$. Indicando con $(x_{0},y_{0})$ il punto e con
$(x,y)$ il punto variabile, il piano
\begin{equation}\label{eq:piano-tangente}
  z=f(x_{0},y_{0})+f_{x}(x_{0},y_{0})(x-x_{0})+f_{y}(x_{0},y_{0})(y-y_{0})
\end{equation}
si chiama \emph{piano tangente} al grafico di $f$ nel punto $(x_{0},y_{0},f(x_{0},y_{0}))$.
\end{definizione}

La \eqref{eq:differenziabile} è la generalizzazione della
\cref{prop:approssimazione-lineare}: dice che
$f(x+h,y+k)=f(x,y)+f_{x}h+f_{y}k+o\bigl(\sqrt{h^{2}+k^{2}}\bigr)$, cioè che vicino al punto il
grafico di $f$ si confonde con il piano tangente, con un errore trascurabile rispetto alla
distanza. Si può dare una definizione apparentemente più generale, senza richiedere a
priori l'esistenza delle derivate parziali, chiedendo che esistano due numeri $a$, $b$
con $f(x+h,y+k)=f(x,y)+ah+bk+o\bigl(\sqrt{h^{2}+k^{2}}\bigr)$; ponendo $k=0$ si vede subito che
necessariamente $a=f_{x}(x,y)$, e analogamente $b=f_{y}(x,y)$.

\begin{teorema}[La differenziabilità implica la continuità]\label{teo:differenziabile-continua}
Ogni funzione differenziabile in un punto $(x,y)$ è continua in $(x,y)$.
\end{teorema}

\begin{schema}
\Ipotesi{$f$ differenziabile in $(x,y)$.}
\Tesi{$f(x+h,y+k)\to f(x,y)$ per $(h,k)\to(0,0)$.}
\end{schema}

\begin{strategia}
Si procede come per il \cref{teo:derivabile-continua}: l'incremento $f(x+h,y+k)-f(x,y)$ si
scrive come la somma del resto della \eqref{eq:differenziabile}, moltiplicato per
$\sqrt{h^{2}+k^{2}}$, e del differenziale $f_{x}h+f_{y}k$. Entrambi tendono a zero.
\end{strategia}

\begin{dimostrazione}
Poniamo $\rho=\sqrt{h^{2}+k^{2}}$ e, per $(h,k)\neq(0,0)$,
\[
  \varepsilon(h,k)=\frac{f(x+h,y+k)-f(x,y)-f_{x}(x,y)\,h-f_{y}(x,y)\,k}{\rho},
\]
che tende a $0$ per l'ipotesi di differenziabilità. Allora
\[
  \abs{f(x+h,y+k)-f(x,y)}=\Bigl|\varepsilon(h,k)+f_{x}\frac h\rho+f_{y}\frac k\rho\Bigr|\,\rho
  \le\bigl\{\abs{\varepsilon(h,k)}+\abs{f_{x}}+\abs{f_{y}}\bigr\}\,\rho,
\]
perché $\abs h\le\rho$ e $\abs k\le\rho$. Per $(h,k)\to(0,0)$ la quantità tra parentesi
graffe è limitata e $\rho\to0$: il secondo membro tende a zero, e quindi
$f(x+h,y+k)\to f(x,y)$, che è la continuità di $f$ nel punto $(x,y)$.
\end{dimostrazione}

La funzione \eqref{eq:f-discontinua} non è continua nell'origine, e quindi, per il
teorema appena dimostrato, non vi è differenziabile, benché abbia le derivate parziali.

\subsection{Derivate direzionali e significato del gradiente}

\begin{teorema}[Derivate direzionali delle funzioni differenziabili]\label{teo:derivata-direzionale}
\index{derivata!direzionale}
Se $f$ è differenziabile nel punto $P=(x,y)$, allora per ogni versore
$v=(\cos\vartheta,\sen\vartheta)$ esiste la derivata direzionale \eqref{eq:derivata-direzionale},
e risulta
\begin{equation}\label{eq:direzionale-gradiente}
  \frac{\partial f}{\partial v}(P)=f_{x}(P)\cos\vartheta+f_{y}(P)\sen\vartheta .
\end{equation}
Di conseguenza, se $\nabla f(P)\neq(0,0)$, la derivata direzionale è massima, e uguale a
$\abs{\nabla f(P)}$, nella direzione del gradiente, ed è nulla nelle direzioni
perpendicolari al gradiente.
\end{teorema}

\begin{schema}
\Ipotesi{$f$ differenziabile in $P$; $v=(\cos\vartheta,\sen\vartheta)$.}
\Tesi{$\frac{\partial f}{\partial v}(P)=\nabla f(P)\cdot v$; massima nella direzione di
$\nabla f(P)$.}
\end{schema}

\begin{strategia}
Basta scegliere, nella \eqref{eq:differenziabile}, l'incremento $(h,k)=tv$, che ha
lunghezza $\abs t$: il rapporto incrementale lungo la direzione $v$ differisce da
$f_{x}\cos\vartheta+f_{y}\sen\vartheta$ per una quantità infinitesima. La seconda parte è poi
la stessa trigonometria della \cref{prop:gradiente-lineare}.
\end{strategia}

\begin{dimostrazione}
Con $(h,k)=(t\cos\vartheta,t\sen\vartheta)$, $t\neq0$, si ha $\sqrt{h^{2}+k^{2}}=\abs t$ e
\begin{multline*}
  \frac{f(x+t\cos\vartheta,y+t\sen\vartheta)-f(x,y)}{t}-\bigl(f_{x}\cos\vartheta+f_{y}\sen\vartheta\bigr)\\
  =\frac{\abs t}{t}\,\varepsilon(t\cos\vartheta,t\sen\vartheta),
\end{multline*}
dove $\varepsilon$ è la funzione della dimostrazione precedente. Il punto
$(t\cos\vartheta,t\sen\vartheta)$ dista $\abs t$ dall'origine; quindi, per la definizione di
limite \eqref{eq:limite-piano-simboli} applicata alla \eqref{eq:differenziabile}, per ogni
$\eta>0$ esiste $\delta>0$ tale che $\abs{\varepsilon(t\cos\vartheta,t\sen\vartheta)}<\eta$ non appena
$0<\abs t<\delta$, cioè $\varepsilon(t\cos\vartheta,t\sen\vartheta)\to0$ per $t\to0$; mentre $\frac{\abs t}t=\pm1$ è limitata: il secondo membro tende a zero, e si ottiene la
\eqref{eq:direzionale-gradiente}. Per la seconda parte, scritto
$\nabla f(P)=r(\cos\varphi,\sen\varphi)$ con $r=\abs{\nabla f(P)}>0$, si ha come nella
\cref{prop:gradiente-lineare} $\frac{\partial f}{\partial v}=r\cos(\vartheta-\varphi)$, che è
massima per $\vartheta=\varphi$ e nulla per $\vartheta=\varphi\pm\frac\pi2$.
\end{dimostrazione}

Le direzioni perpendicolari al gradiente sono quelle in cui $f$, al primo ordine, non
varia: sono le direzioni delle tangenti alle curve di livello. Il gradiente è quindi
perpendicolare alle curve di livello e punta verso i valori crescenti di $f$, come si
vede nella \cref{fig:livello} (dove sono disegnati i vettori $-\nabla p$, che puntano
invece verso le pressioni decrescenti). Per la funzione \eqref{eq:f-discontinua}, che non è
differenziabile nell'origine, la formula \eqref{eq:direzionale-gradiente} è falsa: si
avrebbe $\frac{\partial f}{\partial v}(0,0)=0$ in ogni direzione, mentre, ad esempio, lungo
la bisettrice il rapporto incrementale $\frac{f(t/\sqrt2,t/\sqrt2)}{t}=\frac1{2t}$ non ha
limite.

\begin{attenzione}[Le derivate direzionali non bastano]
Il \cref{teo:derivata-direzionale} non si inverte: una funzione può avere in un punto la
derivata direzionale secondo \emph{ogni} direzione senza essere differenziabile, e
nemmeno continua. Riprendiamo la funzione $g(x,y)=\frac{x^{2}y}{x^{4}+y^{2}}$
dell'avvertenza «Le rette non bastano», ponendo $g(0,0)=0$. Per
$v=(\cos\vartheta,\sen\vartheta)$ con $\sen\vartheta\neq0$ si ha
\[
  \frac{g(t\cos\vartheta,t\sen\vartheta)}{t}
  =\frac{\cos^{2}\vartheta\,\sen\vartheta}{t^{2}\cos^{4}\vartheta+\sen^{2}\vartheta},
\]
che tende a $\frac{\cos^{2}\vartheta}{\sen\vartheta}$ per $t\to0$; per $\sen\vartheta=0$
la funzione è nulla e la derivata direzionale vale $0$. Tutte le derivate direzionali
esistono, ma $g$ non è continua nell'origine (vale $\frac12$ lungo la parabola $y=x^{2}$),
e quindi non è differenziabile; coerentemente, la \eqref{eq:direzionale-gradiente} è
falsa, perché $\nabla g(0,0)=(0,0)$. Non basta neppure che la
\eqref{eq:direzionale-gradiente} sia verificata: per la funzione
$\frac{x^{4}y}{x^{8}+y^{2}}$, posta uguale a $0$ nell'origine, tutte le derivate
direzionali nell'origine sono nulle, come il gradiente, eppure lungo la curva $y=x^{4}$
la funzione vale costantemente $\frac12$. Per stabilire la differenziabilità serve
dunque la definizione, oppure il teorema del differenziale che segue.
\end{attenzione}

\subsection{Il teorema del differenziale}

Verificare la differenziabilità con la definizione è spesso laborioso. Un importante
criterio è espresso dal teorema seguente.

\begin{teorema}[del differenziale]\label{teo:differenziale}\index{teorema!del differenziale}
Se $f$ ammette derivate parziali prime in un intorno del punto $(x,y)$, e queste sono
continue in $(x,y)$, allora $f$ è differenziabile in $(x,y)$.
\end{teorema}

\begin{schema}
\Ipotesi{$f_{x}$, $f_{y}$ esistono in un intorno di $(x,y)$ e sono continue in $(x,y)$.}
\Tesi{vale la \eqref{eq:differenziabile}.}
\end{schema}

\begin{strategia}
Andiamo dal punto $(x,y)$ al punto $(x+h,y+k)$ in due tappe, parallele agli assi:
prima da $(x,y)$ a $(x,y+k)$ in verticale, poi da $(x,y+k)$ a $(x+h,y+k)$ in
orizzontale. Su ciascuna tappa la funzione dipende da una sola variabile, e il teorema di
Lagrange esprime l'incremento con una derivata parziale calcolata in un punto
intermedio. La continuità delle derivate parziali garantisce che queste differiscano di
poco da quelle nel punto $(x,y)$.
\end{strategia}

\begin{dimostrazione}
Sia $I_{\delta}$ un intorno di $(x,y)$ in cui esistono $f_{x}$ e $f_{y}$, e sia
$\sqrt{h^{2}+k^{2}}<\delta$. Applicando due volte il teorema di Lagrange (relativo a funzioni
di una variabile) abbiamo
\begin{equation}\label{eq:differenziale-lagrange}
\begin{split}
  f(x+h,y+k)-f(x,y)&=\bigl[f(x+h,y+k)-f(x,y+k)\bigr]\\
  &\quad+\bigl[f(x,y+k)-f(x,y)\bigr]\\
  &=f_{x}(x(h),y+k)\,h+f_{y}(x,y(k))\,k,
\end{split}
\end{equation}
dove $x(h)$ è un punto opportuno dell'intervallo di estremi $x$, $x+h$, mentre $y(k)$ è un
punto opportuno dell'intervallo di estremi $y$, $y+k$ (i segmenti percorsi stanno in
$I_{\delta}$, per lo stesso argomento usato nella dimostrazione del teorema di
Schwarz). Si noti in particolare che $x(h)\to x$ e $y(k)\to y$ per $(h,k)\to(0,0)$. Si
ottiene la stima
\begin{equation}\label{eq:differenziale-stima}
\begin{split}
  &\Bigl|\frac{f(x+h,y+k)-f(x,y)-f_{x}(x,y)h-f_{y}(x,y)k}{\sqrt{h^{2}+k^{2}}}\Bigr|\\
  &\quad\le\abs{f_{x}(x(h),y+k)-f_{x}(x,y)}\frac{\abs h}{\sqrt{h^{2}+k^{2}}}\\
  &\qquad+\abs{f_{y}(x,y(k))-f_{y}(x,y)}\frac{\abs k}{\sqrt{h^{2}+k^{2}}}\\
  &\quad\le\abs{f_{x}(x(h),y+k)-f_{x}(x,y)}+\abs{f_{y}(x,y(k))-f_{y}(x,y)} .
\end{split}
\end{equation}
Per l'ipotesi di continuità delle derivate parziali $f_{x}$ e $f_{y}$ nel punto $(x,y)$,
l'ultimo membro della \eqref{eq:differenziale-stima} tende a zero per $(h,k)\to(0,0)$,
perché i punti $(x(h),y+k)$ e $(x,y(k))$ tendono a $(x,y)$; perciò anche il primo membro
tende a zero.
\end{dimostrazione}

Se una funzione $f$ è continua in un insieme aperto $D$ si dice anche che è \emph{di
classe $C^{0}$ in $D$} e si scrive $f\in C^{0}(D)$; se ammette derivate parziali prime
continue in $D$ si dice che è \emph{di classe $C^{1}$ in $D$} e si scrive $f\in C^{1}(D)$;
più in generale $f\in C^{k}(D)$, con $k\in\N$, significa che $f$ ammette derivate
parziali continue in $D$ fino all'ordine $k$. Infine, se $f\in C^{k}(D)$ per ogni $k$, si
scrive $f\in C^{\infty}(D)$. Con queste notazioni, il teorema del differenziale e la
proprietà di continuità si riassumono così:
\begin{equation}\label{eq:classi-differenziabilita}
  f\in C^{1}(D)\ \implica\ f \text{ differenziabile in } D\ \implica\ f\in C^{0}(D).
\end{equation}
Nessuna delle due implicazioni si inverte. L'\cref{es:non-differenziabile} presenta una
funzione continua non differenziabile. Una funzione differenziabile ovunque ma non di
classe $C^{1}$ è invece $f(x,y)=x^{2}\sen\frac1x$ per $x\neq0$, $f(0,y)=0$: è la funzione $f_{2}$
dell'\cref{oss:f2-non-C1}, pensata come funzione di due variabili. Nei punti $(0,y)$ si ha
$f(h,y+k)-f(0,y)=h^{2}\sen\frac1h$ (o $0$ se $h=0$), che in valore assoluto non supera
$h^{2}\le h^{2}+k^{2}$: con $f_{x}(0,y)=f_{y}(0,y)=0$ la \eqref{eq:differenziabile} è verificata, e
$f$, che fuori dall'asse $y$ è di classe $C^{1}$, è differenziabile ovunque; ma $f_{x}=2x\sen\frac1x-\cos\frac1x$ non è continua nei
punti dell'asse $y$, e quindi $f\notin C^{1}(\R^{2})$.

\begin{esempio}[Propagazione degli errori di misura]\label{es:propagazione-errori}
Il differenziale fornisce la regola pratica per stimare l'errore su una grandezza
calcolata da altre misurate. La densità dell'aria si ricava dalla pressione $p$ e dalla
temperatura assoluta $T$ con l'equazione di stato, $\rho=\frac p{RT}$. Le derivate
parziali sono $\rho_{p}=\frac1{RT}=\frac\rho p$ e $\rho_{T}=-\frac p{RT^{2}}=-\frac\rho T$, continue
per $T>0$; per il teorema del differenziale
\[
  \Delta\rho\approx\rho_{p}\Delta p+\rho_{T}\Delta T, \qquad
  \frac{\Delta\rho}{\rho}\approx\frac{\Delta p}{p}-\frac{\Delta T}{T}, \qquad
  \Bigl|\frac{\Delta\rho}{\rho}\Bigr|\lesssim\frac{\abs{\Delta p}}{p}+\frac{\abs{\Delta T}}{T} :
\]
gli errori relativi, al primo ordine, si sommano. Il programma confronta questa stima
con una simulazione in cui $p$ e $T$ variano a caso nei rispettivi intervalli di
incertezza.
\end{esempio}

\begin{python}[Propagazione degli errori: differenziale e simulazione]
import numpy as np

R = 287.05
rho = lambda p, T: p / (R * T)           # densità dell'aria
p0, T0, dp, dT = 101325.0, 288.15, 300.0, 1.5   # misure e incertezze
r0 = rho(p0, T0)
lin = r0 * (dp / p0 + dT / T0)           # stima con il differenziale
rng = np.random.default_rng(1)
p = p0 + rng.uniform(-dp, dp, 100000)
T = T0 + rng.uniform(-dT, dT, 100000)
print(f"rho = {r0:.5f} kg/m^3")
print(f"errore massimo (differenziale): {lin:.5f}"
      f"  ({lin / r0:.2%})")
print(f"errore massimo osservato (simulazione): "
      f"{np.max(np.abs(rho(p, T) - r0)):.5f}")
\end{python}
\begin{risultato}
rho = 1.22501 kg/m^3
errore massimo (differenziale): 0.01000  (0.82%)
errore massimo osservato (simulazione): 0.01003
\end{risultato}

La stima lineare coincide con l'errore massimo osservato a meno di qualche unità del
quinto decimale: la piccola differenza è dovuta ai termini del second'ordine che il
differenziale trascura (il caso peggiore esatto, con $p=p_{0}+\Delta p$ e $T=T_{0}-\Delta T$,
dà $0{,}010056$).

\subsection{La derivata lungo un segmento e la formula di Taylor}

Per dimostrare il criterio dell'hessiano studiamo $f$ lungo un segmento uscente dal
punto critico: ci servono la derivata prima e seconda di una funzione di due variabili
lungo un segmento, e una formula di Taylor di ordine due con un resto più preciso di
quello di Peano.

\begin{lemma}[Derivate lungo un segmento]\label{lem:derivata-segmento}
Sia $f\in C^{1}\bigl(I_{\delta}(P_{0})\bigr)$, con $P_{0}=(x_{0},y_{0})$, e sia $(h,k)$ con
$0<\sqrt{h^{2}+k^{2}}<\delta$. La funzione di una variabile
\begin{equation}\label{eq:g-segmento}
  g(t)=f(x_{0}+th,\,y_{0}+tk)
\end{equation}
è derivabile in un intervallo aperto contenente $[0,1]$, con
\begin{equation}\label{eq:g-primo}
  g'(t)=f_{x}(P_{t})\,h+f_{y}(P_{t})\,k, \qquad P_{t}=(x_{0}+th,\,y_{0}+tk).
\end{equation}
Se inoltre $f\in C^{2}\bigl(I_{\delta}(P_{0})\bigr)$, allora $g$ è derivabile due volte, $g''$ è
continua e
\begin{equation}\label{eq:g-secondo}
  g''(t)=f_{xx}(P_{t})\,h^{2}+2f_{xy}(P_{t})\,hk+f_{yy}(P_{t})\,k^{2} .
\end{equation}
\end{lemma}

\begin{schema}
\Ipotesi{$f\in C^{1}$ (risp.\ $C^{2}$) in $I_{\delta}(P_{0})$; $0<\sqrt{h^{2}+k^{2}}<\delta$.}
\Tesi{$g'(t)=f_{x}(P_{t})h+f_{y}(P_{t})k$; \ (se $f\in C^{2}$)
$g''(t)=f_{xx}h^{2}+2f_{xy}hk+f_{yy}k^{2}$ in $P_{t}$, continua.}
\end{schema}

\begin{strategia}
La derivata di $g$ è la derivata di $f$ nella direzione del segmento, moltiplicata per la
lunghezza di $(h,k)$: per le funzioni differenziabili (e le funzioni di classe $C^{1}$ lo
sono, per il teorema del differenziale) essa si calcola con il gradiente. La derivata
seconda si ottiene applicando lo stesso risultato alla funzione $hf_{x}+kf_{y}$, e il
teorema di Schwarz riunisce i due termini misti.
\end{strategia}

\begin{dimostrazione}
Posto $r=\sqrt{h^{2}+k^{2}}$, i punti $P_{t}$ stanno in $I_{\delta}(P_{0})$ per
$\abs t<\frac\delta r$, e $\frac\delta r>1$. Sia $v=\frac1r(h,k)$, che è un versore. Per il
teorema del differenziale $f$ è differenziabile in ogni punto $P_{t}$, e per il
\cref{teo:derivata-direzionale}
\begin{align*}
  g'(t)&=\lim_{s\to0}\frac{f(P_{t}+s(h,k))-f(P_{t})}{s}
  =r\lim_{s\to0}\frac{f(P_{t}+sr\,v)-f(P_{t})}{sr}\\
  &=r\,\frac{\partial f}{\partial v}(P_{t})=f_{x}(P_{t})h+f_{y}(P_{t})k .
\end{align*}
Se $f\in C^{2}$, la funzione $\varphi=hf_{x}+kf_{y}$ è di classe $C^{1}$, e $g'(t)=\varphi(P_{t})$.
Applicando a $\varphi$ quanto appena dimostrato,
$g''(t)=h\,\varphi_{x}(P_{t})+k\,\varphi_{y}(P_{t})=h(hf_{xx}+kf_{yx})+k(hf_{xy}+kf_{yy})$, calcolato in
$P_{t}$, che per il teorema di Schwarz ($f_{xy}=f_{yx}$) è la \eqref{eq:g-secondo}. Infine
$g''$ è continua: se $t_{n}\to t$, allora $P_{t_{n}}\to P_{t}$ (le coordinate convergono,
\eqref{eq:distanza-componenti}), e per la continuità delle derivate seconde
(\cref{oss:limite-piano-successioni}) $f_{xx}(P_{t_{n}})\to f_{xx}(P_{t})$, e così per le
altre; quindi $g''(t_{n})\to g''(t)$.
\end{dimostrazione}

\begin{lemma}[Formula di Taylor con il resto di Lagrange]\label{lem:taylor-lagrange}
\index{formula di Taylor!con il resto di Lagrange}
Sia $g$ derivabile due volte in un intervallo aperto contenente $[0,1]$. Esiste
$\tau\in(0,1)$ tale che
\begin{equation}\label{eq:taylor-lagrange}
  g(1)=g(0)+g'(0)+\frac12g''(\tau) .
\end{equation}
\end{lemma}

\begin{schema}
\Ipotesi{$g$ derivabile due volte in un intervallo aperto contenente $[0,1]$.}
\Tesi{esiste $\tau\in(0,1)$ con $g(1)=g(0)+g'(0)+\frac12g''(\tau)$.}
\end{schema}

\begin{strategia}
È la formula di Taylor di ordine $1$ con il resto scritto in un punto intermedio, come
nella formula degli accrescimenti finiti \eqref{eq:accrescimenti-finiti}, che ne è il
caso di ordine $0$. Il resto $g(1)-g(0)-g'(0)$ si confronta con $t^{2}$ mediante il teorema
di Cauchy, che produce un rapporto tra $g'$ e $2t$; un'ulteriore applicazione del
teorema di Lagrange a $g'$ fa comparire la derivata seconda.
\end{strategia}

\begin{dimostrazione}
Applichiamo il teorema di Cauchy (\cref{teo:cauchy}) nell'intervallo $[0,1]$ alle
funzioni $G(t)=g(t)-g(0)-g'(0)\,t$ e $K(t)=t^{2}$, osservando che $K'(t)=2t\neq0$ per
$t\in(0,1)$. Poiché $G(0)=K(0)=0$ e $K(1)=1$, esiste $\tau_{1}\in(0,1)$ tale che
\[
  g(1)-g(0)-g'(0)=\frac{G(1)-G(0)}{K(1)-K(0)}=\frac{G'(\tau_{1})}{K'(\tau_{1})}
  =\frac{g'(\tau_{1})-g'(0)}{2\tau_{1}} .
\]
Per il teorema di Lagrange applicato a $g'$ nell'intervallo $[0,\tau_{1}]$, esiste
$\tau\in(0,\tau_{1})$ con $g'(\tau_{1})-g'(0)=g''(\tau)\,\tau_{1}$; sostituendo si ottiene la
\eqref{eq:taylor-lagrange}.
\end{dimostrazione}

\subsection{La dimostrazione del criterio dell'hessiano}
\label{sec:dim-hessiano}

Siamo ora in grado di dimostrare il \cref{teo:hessiano}. Indichiamo, per ogni punto $P$
vicino a $(x_{0},y_{0})$ e ogni $(h,k)$, la forma quadratica
\begin{equation}\label{eq:forma-quadratica}
  Q_{P}(h,k)=f_{xx}(P)\,h^{2}+2f_{xy}(P)\,hk+f_{yy}(P)\,k^{2} ,
\end{equation}
che per il \cref{lem:derivata-segmento} è la derivata seconda di $f$ lungo il segmento di
direzione $(h,k)$.

\begin{dimostrazione}[Dimostrazione del \cref{teo:hessiano}]
\emph{Passo 1: segno della forma quadratica.} Siano $A=f_{xx}(P)$, $B=f_{xy}(P)$,
$C=f_{yy}(P)$. Se $A\neq0$, completando il quadrato si ha l'identità
\begin{equation}\label{eq:completamento-quadrato}
\begin{split}
  A\,Q_{P}(h,k)&=A^{2}h^{2}+2ABhk+ACk^{2}\\
  &=(Ah+Bk)^{2}+(AC-B^{2})k^{2} .
\end{split}
\end{equation}
Se $A>0$ e $AC-B^{2}=Hf(P)>0$, il secondo membro è positivo per $(h,k)\neq(0,0)$ (se $k=0$
il primo addendo è $A^{2}h^{2}>0$), e quindi $Q_{P}(h,k)>0$.

\emph{Passo 2: il minimo.} Supponiamo $Hf(x_{0},y_{0})>0$ e $f_{xx}(x_{0},y_{0})>0$. Le
funzioni $Hf$ e $f_{xx}$ sono continue (\cref{teo:continue-piano}); per il teorema della
permanenza del segno, che vale per le funzioni continue di due variabili con la stessa
dimostrazione del caso di una variabile, esiste $\delta>0$ tale che $Hf>0$ e $f_{xx}>0$ in
tutto $I_{\delta}(x_{0},y_{0})$, intorno nel quale supponiamo $f$ di classe $C^{2}$. Sia
$0<\sqrt{h^{2}+k^{2}}<\delta$ e sia $g(t)=f(x_{0}+th,y_{0}+tk)$. Per il
\cref{lem:derivata-segmento}, $g'(0)=f_{x}(x_{0},y_{0})h+f_{y}(x_{0},y_{0})k=0$, perché il
punto è critico; per il \cref{lem:taylor-lagrange} esiste $\tau\in(0,1)$ con
\[
  f(x_{0}+h,y_{0}+k)-f(x_{0},y_{0})=g(1)-g(0)=\frac12g''(\tau)=\frac12Q_{P_{\tau}}(h,k)>0,
\]
dove l'ultima disuguaglianza segue dal passo 1, perché $P_{\tau}\in I_{\delta}$. Dunque
$f(x,y)>f(x_{0},y_{0})$ per ogni $(x,y)\neq(x_{0},y_{0})$ in $I_{\delta}$: il punto è di minimo
relativo stretto.

\emph{Passo 3: il massimo.} Se $Hf(x_{0},y_{0})>0$ e $f_{xx}(x_{0},y_{0})<0$, si applica il
passo 2 alla funzione $-f$, che ha lo stesso hessiano (il prodotto $f_{xx}f_{yy}$ e il
quadrato $f_{xy}^{2}$ non cambiano cambiando segno a $f$) e derivata seconda
$(-f)_{xx}=-f_{xx}>0$.

\emph{Passo 4: la sella.} Sia $Hf(x_{0},y_{0})=AC-B^{2}<0$, con $A$, $B$, $C$ calcolati in
$(x_{0},y_{0})$. Troviamo due direzioni lungo le quali la forma quadratica ha segni
opposti. Se $A\neq0$, si prenda $(h,k)=(1,0)$, che dà $Q=A$, e $(h,k)=(-B,A)$, che per la
\eqref{eq:completamento-quadrato} dà $A\,Q=(AC-B^{2})A^{2}$, cioè $Q=A(AC-B^{2})$, di segno
opposto ad $A$. Se $A=0$, allora $-B^{2}<0$, cioè $B\neq0$, e $Q(1,k)=k(2B+Ck)$ ha, per $k$
piccolo, il segno di $2Bk$, che cambia con il segno di $k$. Poiché
$Q(\lambda h,\lambda k)=\lambda^{2}Q(h,k)$, moltiplicando $(h,k)$ per un $\lambda>0$ piccolo il
segno della forma non cambia: possiamo quindi supporre $\sqrt{h^{2}+k^{2}}<\delta$, dove
$I_{\delta}$ è un intorno in cui $f$ ha derivate prime e seconde continue. Sia dunque $(h,k)$
una tale direzione con $Q_{(x_{0},y_{0})}(h,k)>0$, e sia $g(t)=f(x_{0}+th,y_{0}+tk)$: si ha $g'(0)=0$ e
$g''(0)=Q_{(x_{0},y_{0})}(h,k)>0$, con $g''$ continua, e per il criterio della derivata
seconda (\cref{teo:test-derivata-seconda}) $t=0$ è un punto di minimo relativo stretto
di $g$: lungo questa direzione, arbitrariamente vicino a $(x_{0},y_{0})$, la funzione assume
valori maggiori di $f(x_{0},y_{0})$. Lungo una direzione con $Q<0$, allo stesso modo,
assume valori minori. Dunque $(x_{0},y_{0})$ non è né di massimo né di minimo relativo.
\end{dimostrazione}

\begin{notastorica}[Il differenziale da Leibniz a Fréchet]
Per Leibniz il differenziale $dy$ era un incremento «infinitamente piccolo», e la formula
$dz=\frac{\partial z}{\partial x}dx+\frac{\partial z}{\partial y}dy$ fu usata liberamente per
tutto il Settecento. La distinzione tra derivate parziali e differenziabilità non poteva
allora neppure emergere: le funzioni studiate erano espressioni analitiche «regolari», per
le quali le due nozioni vanno di pari passo. Quando, nell'Ottocento, con la revisione
critica avviata da Cauchy, gli infinitesimi furono sostituiti dai limiti e si cominciò a
considerare funzioni qualsiasi, ci si accorse che l'esistenza delle derivate parziali non
basta a giustificare quella formula, e che occorre una condizione più forte. La definizione moderna di
differenziabilità, che chiede che l'incremento sia approssimato da una funzione lineare
con un errore trascurabile rispetto alla distanza, fu formulata da Otto Stolz nel 1893 e
ripresa all'inizio del Novecento da William Henry Young e da Maurice Fréchet, che la
estese agli spazi di funzioni: è la nozione che sta alla base dell'analisi moderna e dei
metodi numerici, dai metodi di Newton in più variabili ai codici di fluidodinamica
computazionale.
\end{notastorica}

\subsection{Approfondimento: funzioni omogenee e teorema di Eulero}
\label{subsec:funzioni-omogenee}
\index{funzione!omogenea}\index{teorema!di Eulero sulle funzioni omogenee}

Molte formule dell'ingegneria sono prodotti di potenze delle variabili: la portanza
$\frac12\rho SV^{2}C_{L}$, o il coefficiente di resistenza indotta
$C_{Di}=\frac{C_{L}^{2}}{\pi e\,AR}$ dell'\cref{es:sensibilita-AR}. Moltiplicando tutte le
variabili per uno stesso fattore $t>0$, la funzione risulta moltiplicata per una potenza
di $t$. Queste funzioni si dicono \emph{omogenee}, e il loro gradiente soddisfa
un'identità notevole, scoperta da Eulero (si veda la nota storica del \S\ref{sec:schwarz}).
Come nel resto del paragrafo, ci limitiamo alle funzioni di due variabili.

\begin{definizione}[Funzione omogenea]\label{def:funzione-omogenea}
Un insieme $D\subseteq\R^{2}\setminus\{(0,0)\}$ si dice un \emph{cono} se, per ogni
$(x,y)\in D$ e ogni $t>0$, anche $(tx,ty)\in D$. Sia $D$ un cono aperto e
sia $\alpha\in\R$. Una funzione $f:D\to\R$ si dice \emph{omogenea di grado $\alpha$} se
\begin{equation}\label{eq:omogenea}
  f(tx,ty)=t^{\alpha}f(x,y) \qquad\perogni(x,y)\in D,\ \perogni t>0 .
\end{equation}
\end{definizione}

Sono omogenei di grado $n$, nel cono $\R^{2}\setminus\{(0,0)\}$, i polinomi i cui monomi
hanno tutti grado $n$, come $x^{2}+3xy$ ($n=2$). La norma $\sqrt{x^{2}+y^{2}}$ e la funzione
$\sqrt{\abs{xy}}$ dell'\cref{es:non-differenziabile} sono omogenee di grado $1$: i loro
grafici sono coni con il vertice nell'origine. La funzione $\frac{xy}{x^{2}+y^{2}}$
dell'\cref{es:xy-su-x2y2} e l'angolo polare $\operatorname{arctg}\frac yx$, nel semipiano
$x>0$, sono omogenee di grado $0$; e una funzione omogenea di grado $0$ è costante su ogni
semiretta uscente dall'origine, perché $f(tx,ty)=f(x,y)$. Il potenziale della doppietta
$\frac{x}{x^{2}+y^{2}}$ ha grado $-1$; il coefficiente $C_{Di}(C_{L},AR)$, definito nel
semipiano $AR>0$, ha grado $2-1=1$.

\begin{teorema}[di Eulero sulle funzioni omogenee]\label{teo:eulero-omogenee}
Sia $D$ un cono aperto e sia $\alpha\in\R$.
\begin{enumerate}[label=\textup{(\roman*)}]
  \item Se $f:D\to\R$ è omogenea di grado $\alpha$ ed è differenziabile nel punto
  $(x,y)\in D$, allora
  \begin{equation}\label{eq:eulero-omogenee}
    x\,f_{x}(x,y)+y\,f_{y}(x,y)=\alpha\,f(x,y) .
  \end{equation}
  \item Viceversa, se $f\in C^{1}(D)$ e la \eqref{eq:eulero-omogenee} vale in ogni punto di
  $D$, allora $f$ è omogenea di grado $\alpha$.
\end{enumerate}
\end{teorema}

\begin{schema}
\Ipotesi{$D$ cono aperto; \ (i) $f$ omogenea di grado $\alpha$, differenziabile in
$(x,y)$; \ (ii) $f\in C^{1}(D)$ e $xf_{x}+yf_{y}=\alpha f$ in tutto $D$.}
\Tesi{(i) $xf_{x}+yf_{y}=\alpha f$ nel punto $(x,y)$; \ (ii) $f(tx,ty)=t^{\alpha}f(x,y)$ per
ogni $(x,y)\in D$ e ogni $t>0$.}
\end{schema}

\begin{strategia}
Tutto si svolge sulla semiretta che esce dall'origine e passa per $(x,y)$, lungo la quale
si studia la funzione di una variabile $\psi(t)=f(tx,ty)$. La sua derivata è, a meno del
fattore $\sqrt{x^{2}+y^{2}}$, la derivata di $f$ nella direzione radiale, che per le
funzioni differenziabili si calcola con il gradiente (\cref{teo:derivata-direzionale}):
si trova $\psi'(t)=xf_{x}(tx,ty)+yf_{y}(tx,ty)$. L'omogeneità dice invece che
$\psi(t)=t^{\alpha}\psi(1)$. Per (i) basta confrontare le due espressioni di $\psi'(1)$; per
(ii) si mostra che $t^{-\alpha}\psi(t)$ ha derivata nulla, e quindi è costante.
\end{strategia}

\begin{dimostrazione}
Sia $(x,y)\in D$, e siano $r=\sqrt{x^{2}+y^{2}}>0$ e
$v=\frac1r(x,y)$, che è un versore. Per $t>0$ il punto $P_{t}=(tx,ty)$ sta in $D$; poniamo
$\psi(t)=f(P_{t})$. Per $s\neq0$ con $t+s>0$ si ha
\[
  \frac{\psi(t+s)-\psi(t)}{s}=\frac{f(P_{t}+s(x,y))-f(P_{t})}{s}
  =r\,\frac{f(P_{t}+sr\,v)-f(P_{t})}{sr} .
\]
Se $f$ è differenziabile in $P_{t}$, per il \cref{teo:derivata-direzionale} l'ultima
frazione tende, per $s\to0$, alla derivata direzionale
$\frac{\partial f}{\partial v}(P_{t})=f_{x}(P_{t})\frac xr+f_{y}(P_{t})\frac yr$; quindi
$\psi$ è derivabile in $t$ e
\begin{equation}\label{eq:psi-radiale}
  \psi'(t)=x\,f_{x}(tx,ty)+y\,f_{y}(tx,ty) .
\end{equation}

(i) Per l'omogeneità $\psi(t)=t^{\alpha}f(x,y)$, da cui $\psi'(1)=\alpha f(x,y)$. D'altra
parte $f$ è differenziabile in $P_{1}=(x,y)$, e la \eqref{eq:psi-radiale} con $t=1$ dà
$\psi'(1)=xf_{x}(x,y)+yf_{y}(x,y)$: confrontando, si ottiene la tesi.

(ii) Per il teorema del differenziale (\cref{teo:differenziale}) $f$ è differenziabile in
ogni punto di $D$, e la \eqref{eq:psi-radiale} vale per ogni $t>0$. Moltiplicandola per
$t$ e applicando l'ipotesi nel punto $P_{t}\in D$, si ottiene
\[
  t\,\psi'(t)=(tx)\,f_{x}(P_{t})+(ty)\,f_{y}(P_{t})=\alpha f(P_{t})=\alpha\,\psi(t).
\]
La funzione $\chi(t)=t^{-\alpha}\psi(t)$ ha allora, per ogni $t>0$, derivata
$\chi'(t)=t^{-\alpha-1}\bigl[t\,\psi'(t)-\alpha\,\psi(t)\bigr]=0$. Per il
\cref{cor:funzioni-costanti}, applicato nell'intervallo di estremi $1$ e $t$, risulta
$\chi(t)=\chi(1)=f(x,y)$, cioè $f(tx,ty)=t^{\alpha}f(x,y)$.
\end{dimostrazione}

\begin{esempio}[Sensibilità relative e grado di omogeneità]\label{es:omogenee-sensibilita}
Se $f>0$, dividendo la \eqref{eq:eulero-omogenee} per $f$ si ottiene
\[
  \frac{x\,f_{x}}{f}+\frac{y\,f_{y}}{f}=\alpha :
\]
\emph{le sensibilità relative di una funzione omogenea hanno per somma il grado}. Per il
coefficiente di resistenza indotta $C_{Di}=\frac{C_{L}^{2}}{\pi e\,AR}$, omogeneo di grado
$1$, le derivate calcolate nell'\cref{es:sensibilita-AR} danno
\[
  \frac{C_{L}}{C_{Di}}\,\frac{\partial C_{Di}}{\partial C_{L}}=2, \qquad
  \frac{AR}{C_{Di}}\,\frac{\partial C_{Di}}{\partial AR}=-1,
\]
la cui somma è proprio $1$. Con il differenziale,
$\frac{\Delta C_{Di}}{C_{Di}}\approx2\,\frac{\Delta C_{L}}{C_{L}}-\frac{\Delta AR}{AR}$: per volare
con un $C_{L}$ maggiore del $10\%$ senza aumentare il coefficiente di resistenza indotta
occorre un allungamento maggiore di circa il $20\%$. Allo stesso modo la portanza
$L(\rho,V,C_{L})$, omogenea di grado $1+2+1=4$, ha le sensibilità relative $1$, $2$, $1$
trovate nell'\cref{es:portanza-tre}: il teorema si estende infatti alle funzioni di $n$
variabili, nella forma $\sum_{i}x_{i}f_{x_{i}}=\alpha f$.
\end{esempio}

\needspace{8\baselineskip}
\begin{python}[Il teorema di Eulero verificato con sympy]
import sympy as sp

x, y = sp.symbols("x y", positive=True)
CL, AR, e = sp.symbols("C_L AR e", positive=True)
grado = lambda f, u, v: sp.simplify(       # per Eulero è il grado
    (u * sp.diff(f, u) + v * sp.diff(f, v)) / f)
for f in (x**2 + 3*x*y, sp.sqrt(x**2 + y**2), x*y / (x**2 + y**2),
          x / (x**2 + y**2)):
    print(f"f = {str(f):18}  (x f_x + y f_y)/f = {grado(f, x, y)}")
CDi = CL**2 / (sp.pi * e * AR)             # resistenza indotta
print("C_Di = C_L^2/(pi e AR): grado =", grado(CDi, CL, AR))
\end{python}
\begin{risultato}
f = x**2 + 3*x*y        (x f_x + y f_y)/f = 2
f = sqrt(x**2 + y**2)   (x f_x + y f_y)/f = 1
f = x*y/(x**2 + y**2)   (x f_x + y f_y)/f = 0
f = x/(x**2 + y**2)     (x f_x + y f_y)/f = -1
C_Di = C_L^2/(pi e AR): grado = 1
\end{risultato}

Le funzioni omogenee sono anche alla base dell'\emph{analisi dimensionale}: una legge
fisica non dipende dalle unità di misura, ed è quindi omogenea rispetto a ciascuna
grandezza fondamentale; così si ricava, per un'ala di forma e assetto assegnati e
trascurando viscosità e compressibilità, che la portanza è proporzionale a $\rho V^{2}S$.

%% ----------------------------------------------------------------------------
%%  §7.8  Esercizi di riepilogo
%% ----------------------------------------------------------------------------
\section{Esercizi di riepilogo}
\label{sec:esercizi-cap7}

Come nei capitoli precedenti, ogni soluzione è accompagnata da una nota sulla
\emph{scelta logica} che la guida.

\needspace{8\baselineskip}
\subsection*{Domini e limiti}

\begin{esercizio}
Determinare il dominio della funzione $f(x,y)=\log(1-x^{2}-y^{2})+\sqrt{y-x^{2}}$.
\begin{soluzione}
Il logaritmo richiede $1-x^{2}-y^{2}>0$, cioè $x^{2}+y^{2}<1$: l'interno del cerchio di
centro l'origine e raggio $1$, circonferenza esclusa. La radice richiede $y\ge x^{2}$: la
regione che sta sopra la parabola $y=x^{2}$, parabola compresa. Il dominio è
l'intersezione delle due regioni: la parte del cerchio aperto che sta sopra la parabola
(o su di essa). Le due curve si incontrano dove $y=x^{2}$ e $x^{2}+y^{2}=1$, cioè
$y^{2}+y-1=0$, $y=\frac{\sqrt5-1}2$: nei due punti
$\bigl(\pm\sqrt{(\sqrt5-1)/2},\,(\sqrt5-1)/2\bigr)$, che stanno sulla circonferenza e sono
quindi esclusi dal dominio.
\emph{Scelta logica:} quando la funzione è somma di più termini, il dominio è
l'intersezione dei
domini dei singoli termini; ogni condizione si traduce in una regione delimitata da una
curva, e conviene disegnarle tutte prima di intersecarle.
\end{soluzione}
\end{esercizio}

\begin{esercizio}\label{es:limite-xy2}
Stabilire se esiste $\displaystyle\lim_{(x,y)\to(0,0)}\frac{xy^{2}}{x^{2}+y^{4}}$.
\begin{soluzione}
Lungo gli assi $y=0$ e $x=0$ la funzione è nulla; lungo la retta $y=mx$, con $m\neq0$,
$\frac{m^{2}x^{3}}{x^{2}+m^{4}x^{4}}=\frac{m^{2}x}{1+m^{4}x^{2}}\to0$. Lungo tutte le rette il limite
è $0$; ma lungo la parabola $x=y^{2}$ la funzione vale $\frac{y^{4}}{2y^{4}}=\frac12$. Per la
\cref{prop:restrizione-curva} il limite non esiste. \emph{Scelta logica:} il denominatore
mette sullo stesso piano $x^{2}$ e $y^{4}$, cioè $x$ e $y^{2}$: la curva «giusta» per
smascherare la funzione è quella lungo la quale i due termini sono dello stesso ordine,
$x=y^{2}$. È lo stesso fenomeno della funzione $g$ dell'avvertenza «Le rette non bastano».
\end{soluzione}
\end{esercizio}

\begin{esercizio}
Calcolare $\displaystyle\lim_{(x,y)\to(0,0)}\frac{\sen(xy)}{x}$, nel dominio $x\neq0$, e
$\displaystyle\lim_{(x,y)\to(0,0)}\frac{x^{3}+y^{3}}{x^{2}+y^{2}}$.
\begin{soluzione}
Per la disuguaglianza $\abs{\sen t}\le\abs t$ si ha
$\Bigl|\frac{\sen(xy)}{x}\Bigr|\le\frac{\abs{xy}}{\abs x}=\abs y\le\sqrt{x^{2}+y^{2}}$, che tende a $0$:
il primo limite vale $0$. Per il secondo, in coordinate polari,
$\Bigl|\frac{\rho^{3}(\cos^{3}\vartheta+\sen^{3}\vartheta)}{\rho^{2}}\Bigr|\le2\rho$, e per la
\cref{prop:coordinate-polari} il limite vale $0$. \emph{Scelta logica:} in entrambi i
casi si cerca una maggiorazione che dipenda soltanto dalla distanza dall'origine; nel
primo caso la scorciatoia $\frac{\sen(xy)}{xy}\to1$ richiederebbe di trattare a parte i
punti con $y=0$, mentre la disuguaglianza li comprende tutti.
\end{soluzione}
\end{esercizio}

\needspace{8\baselineskip}
\subsection*{Derivate parziali, differenziabilità, piano tangente}

\begin{esercizio}
Verificare che la funzione $f(x,y)=\operatorname{arctg}\frac yx$, per $x>0$, soddisfa
l'equazione $x\,f_{x}+y\,f_{y}=0$, e che $g(x,y)=\sqrt{x^{2}+y^{2}}$ soddisfa
$x\,g_{x}+y\,g_{y}=g$ per $(x,y)\neq(0,0)$. Interpretare i risultati.
\begin{soluzione}
Con la regola della catena,
$f_{x}=\frac1{1+y^{2}/x^{2}}\cdot\bigl(-\frac y{x^{2}}\bigr)=-\frac{y}{x^{2}+y^{2}}$ e
$f_{y}=\frac{x}{x^{2}+y^{2}}$, da cui $xf_{x}+yf_{y}=\frac{-xy+xy}{x^{2}+y^{2}}=0$. Per $g$, dalla
\eqref{eq:norma-derivate}, $xg_{x}+yg_{y}=\frac{x^{2}+y^{2}}g=g$. Per interpretare i
risultati, si osservi che la combinazione $xf_{x}+yf_{y}$ è, a meno del fattore
$\sqrt{x^{2}+y^{2}}$, la derivata direzionale nella direzione radiale
(\cref{teo:derivata-direzionale}: entrambe le funzioni sono di classe $C^{1}$, e quindi
differenziabili, fuori dall'origine). La funzione $f$ è l'angolo polare, costante lungo
le semirette uscenti dall'origine, e la sua derivata radiale è nulla; $g$ è la distanza
dall'origine, che lungo una semiretta cresce proporzionalmente. Sono casi particolari
del teorema di Eulero sulle funzioni omogenee (\cref{teo:eulero-omogenee}): $f$ è
omogenea di grado $0$ e $g$ di grado $1$, e il teorema afferma che ogni funzione $h$
differenziabile e omogenea di grado $\alpha$ verifica $xh_{x}+yh_{y}=\alpha h$.
\emph{Scelta logica:} le derivate si calcolano con la regola della catena, tenendo
costante l'altra variabile; per interpretare la combinazione $xf_{x}+yf_{y}$ conviene
riconoscervi una derivata direzionale, quella lungo il raggio.
\end{soluzione}
\end{esercizio}

\begin{esercizio}
Verificare il teorema di Schwarz per la funzione $f(x,y)=e^{xy}\sen y$.
\begin{soluzione}
Si ha $f_{x}=y\,e^{xy}\sen y$ e $f_{y}=x\,e^{xy}\sen y+e^{xy}\cos y$. Derivando la prima
rispetto a $y$ con la regola del prodotto,
$f_{xy}=e^{xy}\sen y+xy\,e^{xy}\sen y+y\,e^{xy}\cos y$; derivando la seconda rispetto a $x$,
$f_{yx}=e^{xy}\sen y+xy\,e^{xy}\sen y+y\,e^{xy}\cos y$. Le due espressioni coincidono, come
previsto dal teorema, perché $f$ è di classe $C^{\infty}$ su $\R^{2}$.
\emph{Scelta logica:} il calcolo delle due derivate miste per strade diverse è un
ottimo controllo
degli errori di calcolo, proprio perché il risultato è noto in anticipo.
\end{soluzione}
\end{esercizio}

\begin{esercizio}\label{es:non-differenziabile}
Mostrare che la funzione $f(x,y)=\sqrt{\abs{xy}}$ è continua su $\R^{2}$ e ha le derivate
parziali nell'origine, ma non è differenziabile nell'origine.
\begin{soluzione}
La continuità segue dal \cref{teo:continue-piano}, perché $f$ è composizione delle
funzioni continue $xy$, $\abs t$, $\sqrt t$. La funzione è nulla sugli assi, quindi
$f_{x}(0,0)=f_{y}(0,0)=0$. Se fosse differenziabile, il limite
\eqref{eq:differenziabile} sarebbe $\lim\frac{\sqrt{\abs{hk}}}{\sqrt{h^{2}+k^{2}}}=0$; ma lungo la
bisettrice $k=h$ il rapporto vale $\frac{\abs h}{\sqrt2\,\abs h}=\frac1{\sqrt2}$, e il limite, se
esistesse, dovrebbe valere $\frac1{\sqrt2}$ (\cref{prop:restrizione-curva}). Dunque $f$ non
è differenziabile in $(0,0)$. L'esempio mostra che la seconda implicazione della
\eqref{eq:classi-differenziabilita} non si inverte. Poiché $f(tx,ty)=\abs t\,f(x,y)$, il
grafico di $f$ è un cono con il vertice nell'origine, e, come $\abs x$ in una variabile,
nel vertice non ammette piano tangente.
\emph{Scelta logica:} per escludere la differenziabilità basta trovare una curva lungo
la quale il rapporto della \eqref{eq:differenziabile} non tende a zero.
\end{soluzione}
\end{esercizio}

\begin{esercizio}
Scrivere l'equazione del piano tangente al paraboloide $z=x^{2}+y^{2}$ nel punto
$(1,1,2)$, e verificare che il paraboloide sta tutto sopra di esso.
\begin{soluzione}
Si ha $f_{x}=2x$, $f_{y}=2y$, quindi $f_{x}(1,1)=f_{y}(1,1)=2$, e la
\eqref{eq:piano-tangente} dà $z=2+2(x-1)+2(y-1)=2x+2y-2$. La differenza tra la funzione e
il piano è $x^{2}+y^{2}-2x-2y+2=(x-1)^{2}+(y-1)^{2}\ge0$, nulla soltanto nel punto di
tangenza. \emph{Scelta logica:} la funzione è un polinomio, quindi di classe $C^{1}$ e
differenziabile ovunque (\cref{teo:differenziale}); la verifica finale è la versione in
due variabili della convessità del \cref{sec:convessita}: il grafico sta sopra il piano
tangente.
\end{soluzione}
\end{esercizio}

\needspace{8\baselineskip}
\subsection*{Estremi e ottimizzazione}

\begin{esercizio}
Determinare e classificare i punti critici di $f(x,y)=x^{3}+y^{3}-3xy$.
\begin{soluzione}
Il sistema $f_{x}=3x^{2}-3y=0$, $f_{y}=3y^{2}-3x=0$ dà $y=x^{2}$ e $x^{4}=x$, cioè $x=0$ oppure
$x=1$: i punti critici sono $(0,0)$ e $(1,1)$. Si ha $f_{xx}=6x$, $f_{xy}=-3$, $f_{yy}=6y$,
quindi $H=36xy-9$. In $(0,0)$, $H=-9<0$: sella. In $(1,1)$, $H=27>0$ ed $f_{xx}=6>0$:
minimo relativo, con $f(1,1)=-1$; non è assoluto, perché $f(x,0)=x^{3}\to-\infty$ per
$x\to-\infty$. I risultati coincidono con l'ultima parte dell'output del programma del
\cref{sec:estremi-due}. \emph{Scelta logica:} nei sistemi polinomiali conviene ricavare
una variabile dall'equazione più semplice e sostituirla nell'altra, riducendosi a
un'equazione in una sola incognita.
\end{soluzione}
\end{esercizio}

\begin{esercizio}[Il serbatoio di superficie minima]\label{es:serbatoio}
Un serbatoio aperto in alto, a forma di parallelepipedo, deve contenere il volume $V$.
Determinare le dimensioni che rendono minima la superficie (e quindi il peso della
lamiera).
\begin{soluzione}
Siano $x$, $y$ i lati di base e $z$ l'altezza, con $xyz=V$. La superficie del fondo e delle
quattro pareti è $xy+2xz+2yz$; ricavando $z=\frac V{xy}$ si ottiene una funzione di due
variabili nel quadrante $x>0$, $y>0$:
\[
  S(x,y)=xy+\frac{2V}y+\frac{2V}x .
\]
I punti critici risolvono $S_{x}=y-\frac{2V}{x^{2}}=0$ e $S_{y}=x-\frac{2V}{y^{2}}=0$; da
$x^{2}y=2V=xy^{2}$ segue $x=y$, e quindi l'unico punto critico è $(x^{*},y^{*})$, con
$x^{*}=y^{*}=\sqrt[3]{2V}$ e $z=\frac V{(x^{*})^{2}}=\frac12\sqrt[3]{2V}$: l'altezza è metà
del lato di base. In tale punto $S_{xx}=\frac{4V}{x^{3}}=2$, $S_{yy}=2$,
$S_{xy}=1$, quindi $H=3>0$ ed $S_{xx}>0$: è un minimo relativo. È anche il minimo
assoluto. Infatti, posto $m=S(x^{*},y^{*})$, $a=\frac{2V}m$ e $b=\frac ma$, si ha
$S(x,y)>m$ se $x\le a$ (perché $S>\frac{2V}x\ge\frac{2V}a=m$), e analogamente se $y\le a$;
e anche se $x,y\ge a$ e uno dei due è $\ge b$ (perché $S>xy\ge ab=m$). Dunque $S>m$ sul
bordo del quadrato $[a,b]\times[a,b]$ e fuori di esso; in particolare $(x^{*},y^{*})$ è
interno al quadrato. Sul quadrato $S$ ha minimo assoluto per il teorema di Weierstrass
(\cref{oss:weierstrass-piano}); questo minimo è $\le m$, quindi non si trova sul bordo, è
interno, ed è dunque un punto critico (\cref{teo:fermat-due}): l'unico, $(x^{*},y^{*})$.
Poiché fuori dal quadrato $S>m$, questo è il minimo assoluto in tutto il quadrante.
\emph{Scelta logica:} il vincolo $xyz=V$ permette di eliminare una variabile e di
ridurre il problema a una funzione di due variabili senza vincoli; il problema con il
vincolo si potrebbe trattare direttamente con il metodo dei moltiplicatori di Lagrange,
che si studia in Analisi Matematica 2.
\end{soluzione}
\end{esercizio}

\begin{esercizio}[La velocità sulla superficie di un cilindro]
Il potenziale della corrente attorno a un cilindro di raggio $1$, investito da una
corrente uniforme di velocità unitaria, è $\varphi(x,y)=x+\frac x{x^{2}+y^{2}}$. Calcolare le
componenti della velocità $u=\varphi_{x}$, $v=\varphi_{y}$, verificare che sulla circonferenza
$x^{2}+y^{2}=1$ la velocità è tangente alla circonferenza, e calcolarne il modulo nel
punto $(0,1)$.
\begin{soluzione}
Si ha $u=1+\frac{y^{2}-x^{2}}{(x^{2}+y^{2})^{2}}$ e $v=-\frac{2xy}{(x^{2}+y^{2})^{2}}$. Sulla
circonferenza unitaria $u=1+y^{2}-x^{2}$ e $v=-2xy$; la componente radiale della velocità
è $ux+vy=x+xy^{2}-x^{3}-2xy^{2}=x(1-x^{2}-y^{2})=0$: la velocità è tangente, cioè la
circonferenza $x^{2}+y^{2}=1$, contorno del cilindro, è una linea di corrente, che il
fluido non attraversa. Nel punto $(0,1)$ si ha
$u=2$, $v=0$: la velocità è il \emph{doppio} di quella della corrente indisturbata. Per il
teorema di Bernoulli, lì la pressione è minima: è il «picco di depressione» che, su un
profilo alare, genera gran parte della portanza. \emph{Scelta logica:} la tangenza si
verifica annullando il prodotto scalare tra la velocità e il vettore $(x,y)$, che è
normale alla circonferenza (è, a meno di un fattore, il gradiente di $x^{2}+y^{2}$, normale
alle curve di livello). Che $\varphi$ sia armonica è stato verificato nel
\cref{sec:schwarz}.
\end{soluzione}
\end{esercizio}

%% ----------------------------------------------------------------------------
%%  §7.9  Punti fondamentali del capitolo
%% ----------------------------------------------------------------------------
\section{Punti fondamentali del capitolo}
\label{sec:sintesi-cap07}

L'idea centrale del capitolo è che il calcolo in più variabili si costruisce riducendosi
alle funzioni di una variabile: si fissa una variabile, o ci si muove lungo una retta o
una curva, e si usano i teoremi già noti. La riduzione però non basta sempre: le
informazioni raccolte lungo poche direzioni non controllano un intero intorno, e la
nozione che lo fa davvero è la differenziabilità.

\begin{puntichiave}
\item \textbf{Funzioni di due variabili.} Una funzione $f:D\to\R$, con
$D\subseteq\R^{2}$, ha per grafico una superficie e si descrive bene con le curve di
livello $\set{(x,y)\in D \tc f(x,y)=c}$; il dominio naturale si trova risolvendo
disequazioni in due variabili (\cref{def:funzione-due-variabili,def:curve-livello}).
\item \textbf{Limiti e continuità.} La definizione con $\varepsilon$ e $\delta$ usa la
distanza nel piano ed equivale a quella con le successioni di punti; polinomi, somme,
prodotti, quozienti e composizioni di funzioni continue sono continui
(\cref{def:limite-piano,oss:limite-piano-successioni,teo:continue-piano}).
\item \textbf{Esistenza e non esistenza del limite.} Se lungo due curve che tendono a
$P_{0}$ la funzione ha limiti diversi, il limite non esiste; ma l'accordo lungo tutte le
rette non basta, come mostra $\frac{x^{2}y}{x^{4}+y^{2}}$ lungo $y=x^{2}$. Per provare
l'esistenza serve una maggiorazione di $\abs{f-\ell}$ che dipenda soltanto dalla
distanza $\rho$ da $P_{0}$ e tenda a $0$
(\cref{prop:restrizione-curva,prop:coordinate-polari}).
\item \textbf{Derivate parziali e gradiente.} $f_{x}$ è la derivata della funzione di
una variabile $x\mapsto f(x,y)$, con $y$ fissato, e si calcola con le regole usuali;
$\nabla f=(f_{x},f_{y})$. L'esistenza delle derivate parziali in un punto non implica la
continuità: $\frac{xy}{x^{2}+y^{2}}$, nulla nell'origine, le ha ma non è continua
(\cref{def:derivate-parziali,def:gradiente} e \eqref{eq:f-discontinua}).
\item \textbf{Teorema di Schwarz.} Se $f_{x}$, $f_{y}$, $f_{xy}$, $f_{yx}$ esistono in un
intorno di $(x_{0},y_{0})$ e le derivate miste sono continue in $(x_{0},y_{0})$, allora
$f_{xy}(x_{0},y_{0})=f_{yx}(x_{0},y_{0})$; senza la continuità le derivate miste possono
differire (\cref{teo:schwarz,oss:schwarz-controesempio}).
\item \textbf{Condizione necessaria per gli estremi.} In un punto di massimo o di minimo
relativo \emph{interno} al dominio, in cui esistono le derivate parziali, il gradiente
è nullo. La condizione non è sufficiente: per $y^{2}-x^{2}$ l'origine è un punto di
sella (\cref{teo:fermat-due}).
\item \textbf{Criterio dell'hessiano.} Se $f$ ha derivate prime e seconde continue
vicino al punto critico e $H=f_{xx}f_{yy}-f_{xy}^{2}$: per $H>0$ e $f_{xx}>0$ si ha un
minimo relativo stretto, per $H>0$ e $f_{xx}<0$ un massimo relativo stretto, per $H<0$
una sella; per $H=0$ il criterio non decide (\cref{teo:hessiano}).
\item \textbf{Minimi quadrati.} La retta $C_{D}=b+aX$ che minimizza la somma dei
quadrati dei residui si ottiene annullando il gradiente di $S(a,b)$: le equazioni
normali \eqref{eq:equazioni-normali} hanno un'unica soluzione se gli $X_{i}$ non sono
tutti uguali, e la soluzione è il minimo assoluto (\cref{fig:minimi-quadrati}).
\item \textbf{Funzioni di $n$ variabili.} Derivate parziali e gradiente si definiscono
allo stesso modo; la matrice hessiana è simmetrica se le derivate seconde sono continue,
e un punto critico è di minimo (massimo) relativo se l'hessiana vi è definita
positiva (negativa) (\cref{def:parziali-n,oss:hessiana-n}).
\item \textbf{Differenziabilità.} $f$ è differenziabile in un punto se l'incremento è
approssimato dal differenziale $f_{x}h+f_{y}k$ con un errore $o\bigl(\sqrt{h^{2}+k^{2}}\bigr)$;
allora $f$ è continua e il suo grafico ha piano tangente
(\cref{def:differenziabile,teo:differenziabile-continua}).
\item \textbf{Derivate direzionali.} Se $f$ è differenziabile, la derivata nella
direzione del versore $v$ è $\nabla f\cdot v$: il gradiente indica la direzione di
massima crescita ed è perpendicolare alle curve di livello. Il viceversa è falso: tutte
le derivate direzionali possono esistere senza che $f$ sia continua
(\cref{teo:derivata-direzionale}).
\item \textbf{Teorema del differenziale.} Se $f_{x}$, $f_{y}$ esistono vicino a un
punto e sono continue nel punto, $f$ vi è differenziabile; dunque $C^{1}$ implica
differenziabile, che implica $C^{0}$, e nessuna delle due implicazioni si inverte
(\cref{teo:differenziale,es:non-differenziabile}).
\item \textbf{Errori da evitare.} Dedurre l'esistenza di un limite dai soli percorsi
rettilinei; scambiare l'esistenza delle derivate parziali o direzionali con la
differenziabilità; applicare il criterio dell'hessiano quando $H=0$, o a un punto non
interno al dominio.
\item \textbf{Applicazioni aeronautiche.} Le derivate parziali misurano la sensibilità
di $C_{D}$ a $C_{L}$ e ad $AR$ e quella della portanza a $\rho$, $V$, $C_{L}$ (in
termini relativi $1$, $2$, $1$); il differenziale stima la propagazione degli errori in
$\rho=\frac p{RT}$, dove gli errori relativi si sommano
(\cref{es:sensibilita-AR,es:portanza-tre,es:propagazione-errori}).
\end{puntichiave}

\paragraph{Il filo del discorso.} Restano aperti, e appartengono a un secondo corso,
gli estremi vincolati (i moltiplicatori di Lagrange), la formula di Taylor in più
variabili e l'integrazione su domini del piano. Il prossimo capitolo torna alle funzioni
di una variabile e affronta il problema inverso della derivazione: come si
\emph{accumula} una grandezza che varia. Le somme inferiori e superiori con cui vi si
costruisce l'integrale definito sono anche il modello degli integrali doppi, con cui si
calcola il volume sotto il grafico di una funzione di due variabili.

